\documentclass[fleqn,10pt]{wlscirep}
\usepackage[utf8]{inputenc}
\usepackage[T1]{fontenc}
\usepackage{amsmath,amssymb}
\usepackage{booktabs}
\usepackage{array}
\usepackage{multirow}
\usepackage{graphicx}
\usepackage{xcolor}
\usepackage{enumitem}
\usepackage{siunitx}
\usepackage{mathtools}
\usepackage{microtype}
\usepackage{hyperref}
\usepackage{algorithm}
\usepackage{algpseudocode}
\usepackage{subcaption}

% --------- local helpers ---------
\newcolumntype{L}[1]{>{\raggedright\arraybackslash}p{#1}}
\newcolumntype{C}[1]{>{\centering\arraybackslash}p{#1}}
\newcolumntype{R}[1]{>{\raggedleft\arraybackslash}p{#1}}

% Symbol/notation shortcuts
\newcommand{\state}{\mathbf{x}}
\newcommand{\obs}{y}
\newcommand{\dt}{\Delta t}
\newcommand{\NIS}{\mathrm{NIS}}
\newcommand{\Reff}{R_{\mathrm{eff}}}
\newcommand{\Radapt}{R_t^{\mathrm{adapt}}}
\newcommand{\nuadapt}{\nu_t^{\mathrm{adapt}}}
\newcommand{\lambdat}{\lambda_t}
\newcommand{\nuparam}{\nu}
\newcommand{\kurtosis}{\hat{\kappa}_t}
\newcommand{\innovseq}{\{e_\tau\}}

\title{NAROSSM: Noise-Adaptive Robust Oscillatory State-Space Filtering for Video-Based Respiratory Trajectory Reconstruction}

\author[1]{Young-Jae Cho}
\author[2,*]{Dae-Yeol Kim}
\affil[1]{Master's Program in AI Computer Engineering, Department of Electrical, Electronic and Computer Engineering, University of Ulsan, Ulsan, Republic of Korea}
\affil[2]{School of ICT Convergence, University of Ulsan, Ulsan, Republic of Korea}

\affil[*]{daeyeol@ulsan.ac.kr}

\keywords{non-contact respiration, oscillatory state-space model, adaptive Kalman filter, Student-$t$ robust update, inference-to-noise alignment, respiratory trajectory reconstruction}

\begin{abstract}
Video-based contactless respiratory monitoring extracts chest-motion proxy
signals that are inherently corrupted by occlusion, tracking drift, burst
motion, and projection-induced harmonic distortion.
These artifacts produce observation noise that is simultaneously
heavy-tailed, heteroscedastic, and non-stationary---properties that
fundamentally violate the Gaussian, time-invariant assumptions underlying
conventional Kalman filters.
We argue that the core challenge in contactless respiratory estimation is not
simply improving rate accuracy, but achieving \emph{inference-to-noise
alignment}: the statistical assumptions of the estimator must match the
actual noise characteristics of the observation signal.
%
We present \emph{Noise-Adaptive Robust Oscillatory State-Space Model}
(NAROSSM), a principled probabilistic framework that pursues this alignment
through five synergistic mechanisms:
(1)~an oscillatory state-space model that represents respiration as a latent
quasi-periodic trajectory with explicit amplitude, phase, and frequency
dynamics;
(2)~innovation-based adaptive measurement noise estimation that
self-calibrates the filter to achieve unit normalised innovation squared
($\NIS \approx 1$), enabling subsequent robust mechanisms to function on
properly calibrated residuals;
(3)~adaptive process noise that uses the estimated measurement noise as a
signal-quality indicator, tightening the filter bandwidth on clean segments
and widening it during noisy intervals;
(4)~a Student-$t$ variational Bayes update with data-driven
degrees-of-freedom adaptation via running kurtosis, which automatically
adjusts the tail heaviness of the assumed noise distribution to match the
empirical innovation statistics;
and (5)~an oscillator-native trust gate that determines observation
reliability from phase coherence and amplitude stability---signals intrinsic
to the state-space model---providing the classical equivalent of learned
temporal attention without the bistable feedback loops that plague
innovation-based gating approaches.
%
A Rauch--Tung--Striebel backward smoother operating on the resulting
\emph{exactly linear} two-dimensional oscillator yields optimal
trajectory reconstruction without linearisation error.
Evaluated on two public benchmarks---COHFACE (160~trials, 40~subjects,
controlled lab) and MAHNOB-HCI (525~trials, 27~subjects,
emotion-elicitation protocol)---across five observation families spanning
optical flow and spatial profile methods,
NAROSSM achieves the lowest frequency-domain MAE
on OF-Farneb\"ack (0.280~BPM vs.\ KF's 0.290) on COHFACE, while
outperforming the Gaussian KF on three of five families in MAE and
four of five in RMSE on MAHNOB-HCI
(e.g., P1D-Quadratic MAE: 5.45 vs.\ 5.50~BPM), with statistically
significant improvements (Wilcoxon $p < 0.05$, Bonferroni-corrected)
on the three Profile1D families.
The adaptive $R_t$ mechanism achieves empirical NIS~$\approx 1$
across all families and both datasets, correctly converging to
dataset-specific noise levels that differ by up to $7\times$ between
benchmarks---eliminating the need for per-dataset measurement noise
tuning.
Cross-dataset $R$ transfer experiments demonstrate that a Gaussian KF
with fixed $R$ calibrated on one dataset becomes miscalibrated on the
other, while NAROSSM adapts automatically.
A controlled outlier injection experiment confirms that NAROSSM's
Student-$t$ update and trust gate provide graceful degradation under
observation corruption, activating protective mechanisms that remain
dormant on clean signals.
Per-frame diagnostic outputs ($\NIS_t$, $\lambdat$, $\Radapt$, $\nuadapt$)
provide a fully auditable inference trail for clinical interpretability.
\end{abstract}

\begin{document}
\flushbottom
\maketitle
\thispagestyle{empty}

% ===================================================================
%  INTRODUCTION
% ===================================================================
\section*{Introduction}

Respiratory monitoring is fundamental to clinical assessment, yet
conventional contact-based methods---nasal thermistors, chest
impedance belts, capnography---impose physical constraints that limit
continuous, unobtrusive monitoring in settings such as sleep studies,
neonatal care, and post-operative surveillance.
Video-based contactless respiratory estimation has emerged as a
promising alternative, leveraging the mechanical coupling between
breathing and visible upper-body motion to extract respiratory
information without physical contact~\cite{Bartula2013,Tarassenko2014,Janssen2016}.

The dominant paradigm treats contactless respiratory estimation as a
\emph{scalar rate extraction problem}: a proxy signal is obtained from
video, bandpass-filtered to the respiratory band, and a single
respiratory rate is estimated via spectral peak detection or
autocorrelation~\cite{Li2014optical,Bai2021}.
However, respiration is not a static scalar quantity.
Each breathing cycle encodes clinically meaningful information in its
\emph{amplitude} (depth of breath), \emph{phase} (timing of
inspiration and expiration), and \emph{frequency dynamics} (gradual or
abrupt changes in breathing rate).
A monitoring system that discards this temporal structure---reducing
continuous physiological dynamics to a single number---sacrifices
both clinical utility and estimation robustness.

We therefore reframe the problem as \emph{oscillatory trajectory
reconstruction}: the goal is to recover the full time-varying
respiratory waveform, including its amplitude, phase, and frequency
modulation, from noisy video-derived observations.

\subsection*{The structural mismatch problem}

This reframing reveals a fundamental challenge.
The proxy signals obtained from video are not direct measurements of
respiratory volume, but rather indirect projections of chest-wall
motion onto a two-dimensional image plane.
Multiple observation extraction methods---dense optical flow,
sparse feature tracking, one-dimensional spatial intensity
profiles---each capture different aspects of this projection with
different noise characteristics~\cite{Farneback2003,LucasKanade1981}.

Empirical analysis of these observation signals reveals noise
properties that are strikingly non-standard:

\begin{itemize}[nosep]
\item \textbf{Heavy tails.}
  Excess kurtosis ranges from approximately~1 (degrees-of-freedom
  profiles) to over~60 (dense optical flow and linear spatial
  profiles), indicating frequent large deviations far exceeding
  Gaussian expectations.
\item \textbf{Heteroscedasticity.}
  ARCH-LM tests reject the null hypothesis of constant variance
  ($p < 0.001$) across all observation families, confirming that
  noise variance is time-varying.
\item \textbf{Non-stationarity.}
  Noise characteristics change qualitatively within a single
  recording due to subject motion, occlusion events, and lighting
  changes.
\end{itemize}

These properties create a \emph{structural mismatch} between the
observation signal and the assumptions of standard temporal estimators.
Conventional Kalman filters assume Gaussian, time-invariant noise;
bandpass-plus-spectral methods assume stationary signal
statistics~\cite{Kalman1960}.
When these assumptions are violated, estimators either become
unstable under heavy-tailed outliers, or become overly conservative
by inflating noise parameters to accommodate worst-case deviations,
thereby ignoring informative observations during normal operation.

\subsection*{The parameter tuning problem}

Classical state-space models encode their distributional assumptions
through a small set of noise parameters---process noise~$Q$,
measurement noise~$R$, and, for robust variants, degrees of
freedom~$\nu$.
These parameters have narrow optimal ranges that vary across
observation methods, recording environments, and subject populations.
For instance, as we demonstrate empirically
(Section~Discussion), the oracle-optimal measurement noise
for a Gaussian Kalman filter differs by over three orders of magnitude
across observation families ($R^* = 0.001$ for optical flow vs.\
$R^* = 1.0$ for spatial profiles on the same dataset) and shifts
substantially between datasets.
No single fixed $R$ can optimise all families and datasets
simultaneously.

In practice, this means that deploying a classical filter on a new
dataset, camera system, or observation method requires a new round
of parameter search---a process that depends on ground-truth
labels and is therefore impractical in clinical settings.
Deep learning approaches circumvent this issue by absorbing
dataset-specific characteristics through data-driven training on
large labelled corpora, but they sacrifice the interpretability and
formal optimality guarantees that make state-space methods
attractive for safety-critical applications.
The challenge, then, is to design a classical state-space framework
whose noise-sensitive parameters \emph{adapt from the data itself},
retaining the transparency of model-based inference while achieving
the environment-invariance of data-driven methods.

\subsection*{Observation-to-physiology alignment}

A prerequisite for addressing the noise mismatch is ensuring that the
observation signal itself is as closely aligned with the physiological
source as possible.
In respiratory monitoring, the physical origin of the observable
signal is the mechanical expansion and contraction of the chest wall
and abdomen during the breathing cycle.

Many video-based approaches derive respiratory information from
indirect sources---facial colour changes, peripheral motion
artefacts---that are only weakly coupled to chest-wall
dynamics~\cite{Tarassenko2014}.
We argue that direct observation of upper-body mechanical motion
provides the strongest physiological basis for respiratory estimation,
and therefore construct our observation pipeline from methods that
explicitly capture chest-wall displacement: dense optical flow
(Farneb\"ack), and one-dimensional spatial intensity profiles at
multiple polynomial orders.

Even with physiologically aligned observation methods, the extracted
signals remain contaminated by the artefact sources described above.
The alignment reduces but does not eliminate the noise mismatch
problem, motivating a principled noise-adaptive inference framework.

\subsection*{Inference-to-noise alignment}

The central contribution of this work is the principle of
\emph{inference-to-noise alignment}: the statistical assumptions
embedded in the estimator should match the actual statistical
properties of the observation noise, and this matching should be
achieved \emph{adaptively} rather than through fixed, pre-specified
parameters.
This directly addresses the parameter tuning problem identified above:
rather than requiring practitioners to find the optimal $R$, $Q$, and
$\nu$ for each new deployment scenario, the framework estimates these
quantities online from the innovation sequence, achieving
near-optimal performance with a single default parameter set across
datasets, frame rates, and observation families.

This principle is operationalised through four concrete mechanisms.
First, respiration is modelled as a latent oscillatory process in a
two-dimensional state space, which is exactly linear given the
instantaneous frequency, eliminating linearisation error from the
filtering equations.
Second, the measurement noise variance is estimated adaptively from
the innovation sequence, ensuring that the filter's internal
calibration reflects the true noise level at each time step.
Third, the degrees of freedom of a Student-$t$ observation likelihood
are adapted via running kurtosis of the innovation sequence, allowing
the estimator to automatically adjust its robustness level: applying
aggressive outlier suppression when innovations are heavy-tailed, and
converging to standard Gaussian filtering when innovations are
well-behaved.
Fourth, an \emph{oscillator-native trust gate} determines
observation reliability using signals intrinsic to the state-space
model---phase coherence and amplitude stability---rather than
innovation magnitude.
In deep learning approaches, temporal attention mechanisms implicitly
learn which observations to trust; in a classical state-space
framework, this capability must be explicitly designed.
Our gate exploits the physical constraint that respiration is
quasi-periodic: observations consistent with the oscillator's
predicted phase progression are trusted, while those causing
abrupt phase jumps or amplitude changes are suppressed.

The resulting framework, which we call NAROSSM (Noise-Adaptive Robust
Oscillatory State-Space Model; Figure~\ref{fig:system_overview}),
makes every noise-related parameter
and trust decision a function of the observed data and the oscillator
state, rather than fixed hyperparameters.
This design ensures that the filter's behaviour is determined by the
actual statistical structure of each observation, not by assumptions
made at design time.

% --- Figure 1: System overview ---
\begin{figure}[!ht]
\centering
% === PLACEHOLDER: Architecture diagram from PaperBanana ===
% \includegraphics[width=\linewidth]{figures/fig1_system_overview}
\fbox{\parbox{0.95\linewidth}{\centering\vspace{4cm}
\textbf{Figure 1: NAROSSM System Overview}\\[6pt]
Video $\to$ Chest Motion Extraction $\to$ Observation Families\\
$\to$ Oscillatory SSM $\to$ Adaptive R + Adaptive $\nu$ + Student-$t$ Update\\
$\to$ RTS Smoother $\to$ Respiratory Trajectory + Diagnostics\\[6pt]
\emph{(Architecture diagram to be created in PaperBanana)}
\vspace{4cm}}}
\caption{%
  Overview of the NAROSSM framework.
  Video frames are processed through physiologically aligned observation
  extraction methods to produce chest-motion proxy signals.
  The proxy signal feeds into an oscillatory state-space model where a
  Kalman filter with innovation-based adaptive $R_t$ and
  kurtosis-adaptive $\nu_t$ performs the forward pass.
  An oscillator-native trust gate ($g_t$) assesses observation
  reliability via phase coherence and amplitude stability, while a
  Student-$t$ variational Bayes update downweights outlier observations
  proportional to their deviation from the calibrated noise model.
  An RTS backward smoother produces the final trajectory estimate.
  Per-frame diagnostic outputs ($\NIS_t$, $\lambdat$, $R_t$, $\nu_t$,
  $g_t$) provide an auditable inference trail.
}
\label{fig:system_overview}
\end{figure}


% ===================================================================
%  RESULTS
% ===================================================================
\section*{Results}

We evaluate NAROSSM on two public benchmarks:
COHFACE~\cite{COHFACE2017} (160~recordings, 40~subjects, 20~fps,
controlled lab setting) and MAHNOB-HCI~\cite{MAHNOB2012}
(525~recordings, 27~subjects, 61~fps, emotion-elicitation protocol).
Five observation extraction families are tested on both datasets: one
optical-flow-based method (OF-Farneb\"ack) and four spatial intensity
profile methods (Profile1D-Linear, -Quadratic, -Cubic, -DoF).
Performance is assessed in both the frequency domain (respiratory rate
accuracy) and the time domain (waveform reconstruction fidelity),
with statistical significance established via paired Wilcoxon
signed-rank tests with Bonferroni correction.
We first present COHFACE results in detail, then MAHNOB-HCI results,
and conclude with a cross-dataset comparison.

\subsection*{Observation noise characterisation}

Before evaluating NAROSSM's performance, we characterise the
statistical properties of the observation noise to establish the
empirical basis for our adaptive design.

% --- Table 1: Noise characterisation ---
\begin{table}[!ht]
\centering
\caption{%
  Statistical characterisation of observation signals across families.
  Excess kurtosis ($\hat{\kappa}$) is shown before and after
  Butterworth bandpass filtering (0.08--0.50~Hz).
  The Jarque--Bera test ($p < 0.05$) rejects Gaussianity for all
  families at both stages.
  Inferred $\hat{\nu}$ shows the Student-$t$ degrees of freedom
  implied by the kurtosis via $\hat{\nu} = 6/\hat{\kappa} + 4$.
  Note the dramatic reduction in kurtosis after bandpass for
  OF-Farneb\"ack and Profile1D families, explaining why the
  NAROSSM adaptive $\nuadapt$ converges to near-Gaussian values.
}
\label{tab:noise_characterisation}
\begin{tabular}{l rr c rr}
\toprule
& \multicolumn{2}{c}{\textbf{Excess kurtosis} $\hat{\kappa}$} &
  & \multicolumn{2}{c}{\textbf{Inferred} $\hat{\nu}$} \\
\cmidrule(lr){2-3} \cmidrule(lr){5-6}
\textbf{Family} &
Pre-BP & Post-BP &
\textbf{JB} $p$ &
Pre-BP & Post-BP \\
\midrule
OF-Farneb\"ack     & 57.2  &  1.3  & ${<}0.001$ &  4.1  &  8.5  \\
Profile1D-Linear   & 69.1  &  2.3  & ${<}0.001$ &  4.1  &  6.6  \\
Profile1D-Quadratic & 13.5  & $-$0.2 & ${<}0.001$ &  4.4  & ${\to}\infty$ \\
Profile1D-Cubic    & 14.3  & $-$0.2 & ${<}0.001$ &  4.4  & ${\to}\infty$ \\
Profile1D-DoF      &  1.0  &  8.2  & ${<}0.001$ & 10.2  &  4.7  \\
\bottomrule
\end{tabular}
\end{table}

Table~\ref{tab:noise_characterisation} reveals a critical interaction
between preprocessing and noise distribution.
Pre-bandpass signals exhibit pronounced non-Gaussianity, with excess
kurtosis reaching 57--69 for optical flow and spatial profile methods.
However, the Butterworth bandpass filter (0.08--0.50~Hz) acts as a
temporal smoothing operator that dramatically reduces kurtosis:
OF-Farneb\"ack drops from $\hat{\kappa} = 57.2$ to $1.3$, and
Profile1D-Quadratic from $13.5$ to $-0.2$ (sub-Gaussian).
The corresponding inferred $\hat{\nu}$ shifts from heavy-tailed
($\hat{\nu} \approx 4$) to near-Gaussian ($\hat{\nu} \to \infty$).

A notable exception is Profile1D-DoF, where kurtosis \emph{increases}
after bandpass (from 1.0 to 8.2), suggesting that the bandpass filter
concentrates energy in modes with structured, non-Gaussian variability.

This finding has a fundamental implication for adaptive filtering
design: when aggressive preprocessing precedes the filter, the
noise distribution seen by the estimator may be substantially
different from the raw observation noise.
NAROSSM's kurtosis-adaptive $\nuadapt$ operates on the
\emph{post-preprocessing} innovation residuals and correctly
identifies the effective (post-bandpass) noise distribution,
converging to near-Gaussian behaviour where appropriate.


\subsection*{Filter calibration validation}

A correctly calibrated filter should produce normalised innovation
squared (NIS) values with mean approximately equal to~1.
We verify this fundamental property for NAROSSM and compare it
against a baseline Gaussian oscillator with fixed measurement noise.

% --- Figure 2: NIS calibration ---
\begin{figure}[!ht]
\centering
% \includegraphics[width=\linewidth]{figures/fig2_nis_calibration}
\fbox{\parbox{0.95\linewidth}{\centering\vspace{5cm}
\textbf{Figure 2: NIS Calibration}\\[6pt]
NAROSSM empirical NIS distributions per observation family.\\
All families achieve NIS mean $\approx 0.95$--$1.0$.\\
Box plots showing per-trial NIS mean and spread.\\
\emph{(To be generated from experimental results)}
\vspace{5cm}}}
\caption{%
  Empirical normalised innovation squared (NIS) distributions for
  NAROSSM across observation families.
  All families achieve NIS mean in the range 0.947--0.995,
  confirming that the innovation-based adaptive $R_t$ correctly
  self-calibrates the filter.
  This calibration property is the prerequisite for the Student-$t$
  mechanism to meaningfully interpret innovation magnitudes: only
  when NIS~$\approx 1$ does a large innovation reliably indicate
  an outlier rather than a miscalibrated noise model.
}
\label{fig:nis_calibration}
\end{figure}


\subsection*{Adaptive noise parameter behaviour}

A key claim of this work is that NAROSSM's noise parameters adapt to
the actual observation statistics without manual tuning.
We examine the temporal behaviour of the adaptive measurement noise
$R_t$ and the adaptive degrees of freedom $\nu_t$
(Figure~\ref{fig:adaptive_params}).

% --- Figure 3: Adaptive R and nu trajectories ---
\begin{figure}[!ht]
\centering
% \includegraphics[width=\linewidth]{figures/fig3_adaptive_params}
\fbox{\parbox{0.95\linewidth}{\centering\vspace{6cm}
\textbf{Figure 3: Adaptive Parameter Trajectories}\\[6pt]
(a) Example trial from OF-Farneb\"ack: $R_t$ time series showing\\
    noise level tracking; $\nu_t$ converging to $\nu_{\max}$ (near-Gaussian)\\[3pt]
(b) Example trial from Profile1D-Quadratic: $R_t$ relatively stable;\\
    $\nu_t$ at $\nu_{\max}$ throughout\\[3pt]
(c) NIS time series for both: empirical NIS $\approx 1.0$,\\
    $\lambda_t \approx 1$ (Student-$t$ inactive on clean data)\\[6pt]
\emph{(To be generated from experimental results)}
\vspace{6cm}}}
\caption{%
  Temporal behaviour of NAROSSM's adaptive noise parameters for
  representative COHFACE trials.
  \textbf{(a)}~OF-Farneb\"ack: adaptive $R_t$ tracks noise level
  changes; the running-kurtosis $\nuadapt$ converges to
  $\nu_{\max} \approx 200$, correctly identifying that post-bandpass
  residuals are near-Gaussian.
  The Student-$t$ weight $\lambdat \approx 1$ throughout, confirming
  that no outlier suppression is applied when the noise is well-behaved.
  \textbf{(b)}~Profile1D-Quadratic: $R_t$ is stable and $\nuadapt$
  reaches $\nu_{\max}$ rapidly, consistent with the low-noise
  characteristics of spatial profile methods.
  \textbf{(c)}~Empirical NIS time series for both families, showing
  calibration around the target value of~1.
  The self-calibrating NIS demonstrates that the adaptive $R_t$
  mechanism functions correctly regardless of the observation
  family's noise level.
}
\label{fig:adaptive_params}
\end{figure}


\subsection*{COHFACE respiratory frequency stability}

Before interpreting performance differences, we characterise a
fundamental property of the COHFACE dataset that strongly influences
baseline behaviour.
We compute the within-trial respiratory frequency standard deviation
across 30-second sliding windows for all 160~trials.

The results reveal striking frequency stability: the median
within-trial frequency standard deviation is 0.000~Hz, and 64\% of
trials (103/160) exhibit \emph{zero} frequency variation across all
evaluation windows.
Even among the remaining 36\% of trials, frequency variability is
modest (mean std = 0.049~Hz, max = 0.063~Hz).

This finding has a critical implication for model comparison.
The Gaussian KF baseline (kfstd) employs a \emph{fixed} transition
frequency $f$ estimated once from the initial spectral peak.
On a dataset where respiratory frequency is essentially constant
within each trial, this fixed-$f$ assumption is not violated, and
the baseline operates under near-ideal conditions.
NAROSSM's time-varying frequency adaptation module, while
architecturally necessary for non-stationary breathing patterns
(e.g., speech, exercise recovery, sleep transitions), provides
no advantage when the underlying frequency is already constant.

% --- Figure: COHFACE frequency stability ---
\begin{figure}[!ht]
\centering
\fbox{\parbox{0.95\linewidth}{\centering\vspace{3cm}
\textbf{Figure: COHFACE Respiratory Frequency Stability}\\[6pt]
(a) Histogram of within-trial frequency std (64\% at zero)\\
(b) Cumulative distribution showing sharp concentration at low variability\\
\emph{(To be generated from experimental results)}
\vspace{3cm}}}
\caption{%
  Within-trial respiratory frequency variability in COHFACE.
  \textbf{(a)}~Histogram of per-trial frequency standard deviation
  across 30-second windows: 64\% of trials exhibit zero variation.
  \textbf{(b)}~Cumulative distribution confirming that COHFACE
  breathing patterns are predominantly stationary, favouring
  fixed-frequency baselines.
}
\label{fig:cohface_freq_stability}
\end{figure}


\subsection*{Frequency-domain performance}

We evaluate respiratory rate estimation accuracy using windowed
mean absolute error (MAE) computed from Welch power spectral density
peak detection over 30-second sliding windows with 1-second stride.

% --- Table 2: Frequency domain results ---
\begin{table}[!ht]
\centering
\caption{%
  Frequency-domain respiratory rate estimation performance (median values).
  MAE = mean absolute error (BPM), RMSE = root mean squared error (BPM),
  $r$ = Pearson correlation, SNR\textsubscript{spec} = spectral
  signal-to-noise ratio (dB), KL = Kullback--Leibler divergence of
  estimated vs.\ ground-truth spectral density.
  Best result per family in \textbf{bold}.
}
\label{tab:freq_performance}
\begin{tabular}{l ccc ccc ccc}
\toprule
& \multicolumn{3}{c}{\textbf{MAE (BPM) $\downarrow$}} &
  \multicolumn{3}{c}{\textbf{RMSE (BPM) $\downarrow$}} &
  \multicolumn{3}{c}{\textbf{Pearson $r$ $\uparrow$}} \\
\cmidrule(lr){2-4} \cmidrule(lr){5-7} \cmidrule(lr){8-10}
\textbf{Family} &
  Base & KF & NAROSSM &
  Base & KF & NAROSSM &
  Base & KF & NAROSSM \\
\midrule
OF-Farneb\"ack &
  0.290 & 0.290 & \textbf{0.280} &
  0.390 & 0.375 & \textbf{0.360} &
  0.890 & 0.900 & \textbf{0.920} \\
P1D-Linear &
  0.335 & 0.315 & \textbf{0.305} &
  0.435 & 0.400 & \textbf{0.395} &
  0.840 & 0.860 & \textbf{0.870} \\
P1D-Quadratic &
  \textbf{0.195} & 0.220 & 0.210 &
  \textbf{0.270} & \textbf{0.270} & \textbf{0.270} &
  \textbf{0.950} & 0.930 & 0.930 \\
P1D-Cubic &
  0.220 & 0.230 & \textbf{0.220} &
  \textbf{0.275} & 0.300 & 0.290 &
  \textbf{0.940} & 0.930 & \textbf{0.940} \\
P1D-DoF &
  2.400 & \textbf{2.105} & 2.180 &
  4.015 & \textbf{3.135} & 3.315 &
  0.370 & \textbf{0.400} & 0.390 \\
\bottomrule
\end{tabular}

\vspace{6pt}

\begin{tabular}{l ccc ccc}
\toprule
& \multicolumn{3}{c}{\textbf{SNR\textsubscript{spec} (dB) $\uparrow$}} &
  \multicolumn{3}{c}{\textbf{KL Div $\downarrow$}} \\
\cmidrule(lr){2-4} \cmidrule(lr){5-7}
\textbf{Family} &
  Base & KF & NAROSSM &
  Base & KF & NAROSSM \\
\midrule
OF-Farneb\"ack &
  5.54 & \textbf{6.91} & 6.74 &
  0.195 & \textbf{0.130} & 0.137 \\
P1D-Linear &
  6.68 & \textbf{8.17} & 8.13 &
  0.180 & \textbf{0.109} & 0.119 \\
P1D-Quadratic &
  7.72 & \textbf{9.27} & 9.15 &
  0.144 & \textbf{0.087} & 0.088 \\
P1D-Cubic &
  7.57 & \textbf{9.18} & 9.06 &
  0.146 & \textbf{0.089} & \textbf{0.089} \\
P1D-DoF &
  2.22 & \textbf{4.33} & 4.20 &
  0.403 & \textbf{0.298} & 0.319 \\
\bottomrule
\end{tabular}
\end{table}

Table~\ref{tab:freq_performance} reveals a family-dependent pattern
consistent with NAROSSM's design philosophy of inference-to-noise alignment.
NAROSSM achieves the best or tied-best frequency MAE on four of five
observation families.
For the noisiest optical-flow family (OF-Farneb\"ack), NAROSSM achieves the
lowest median MAE (0.280~BPM) and RMSE (0.360~BPM), outperforming
both the raw baseline (0.290/0.390) and the Gaussian KF (0.290/0.375),
with the highest Pearson correlation ($r = 0.920$).
This demonstrates that adaptive noise estimation provides genuine
benefit when observation noise is heteroscedastic: NAROSSM's adaptive
$R_t$ correctly tracks the time-varying noise variance, yielding a
better-calibrated Kalman gain than the fixed-$R$ baseline.
On P1D-Linear, NAROSSM further improves over the Gaussian KF
(MAE: 0.305 vs.\ 0.315; $r$: 0.870 vs.\ 0.860), confirming that
adaptive mechanisms benefit intermediate-noise families as well.
For P1D-DoF, the oscillatory SSM with RTS smoothing provides
substantial benefit over the raw signal (MAE from 2.400 to
2.105/2.180), with the Gaussian KF achieving slightly better
performance due to the lower overhead of a fixed noise model on
this highly variable family.

A key finding emerges for high-SNR spatial profile families.
On P1D-Cubic, NAROSSM matches the raw baseline MAE (0.220~BPM)
and Pearson $r$ (0.940),
whereas the Gaussian KF introduces a degradation (MAE: 0.230, $r$: 0.930).
On P1D-Quadratic, NAROSSM outperforms the Gaussian KF in MAE
(0.210 vs.\ 0.220) while all three methods share the same RMSE (0.270).
This ``do no harm'' property arises from NAROSSM's adaptive process noise:
when $R_t$ converges to a low value on clean signals, the adaptive $Q_t$
correspondingly tightens, reducing the filter's temporal smoothing
and preserving the raw signal's spectral accuracy.

The raw spectral baseline retains the lowest MAE on P1D-Quadratic
(0.195~BPM), a structural consequence of Welch frequency quantisation:
with 64\% of COHFACE trials exhibiting zero within-trial frequency
variation, the raw spectral peak is already optimal, and any
temporal filtering introduces bandwidth limitations that shift the
Welch peak by one or two frequency bins.

The spectral quality metrics (SNR\textsubscript{spec} and KL divergence)
reveal an additional nuance: the Gaussian KF achieves marginally higher
spectral SNR and lower KL divergence than NAROSSM on all five families.
This occurs because NAROSSM's time-varying $R_t$ introduces slight
gain fluctuations that modestly broaden the spectral peak relative
to the Gaussian KF's constant gain, which produces a smoother,
more spectrally concentrated output.
Importantly, this spectral smoothness advantage does not translate
to better rate estimation: NAROSSM's superior MAE and Pearson $r$ show
that accurate \emph{peak detection} (which determines the rate estimate)
benefits more from adaptive noise tracking than from spectral purity.
The spectral SNR measures the peak-to-noise ratio of the entire
power spectrum, not the accuracy of peak localisation, and the two
can diverge when adaptive gain yields a slightly broader but more
accurately centred peak.


\subsection*{Time-domain waveform reconstruction}

Waveform reconstruction fidelity is assessed using Lin's concordance
correlation coefficient (CCC)~\cite{LinCCC1989} and time-domain MAE
on bandpass-filtered, cross-correlation-aligned signals.

% --- Table 3: Time domain results ---
\begin{table}[!ht]
\centering
\caption{%
  Time-domain waveform reconstruction performance (median values).
  CCC = Lin's concordance correlation coefficient~\cite{LinCCC1989},
  MAE = mean absolute error, RMSE = root mean squared error,
  SNR\textsubscript{time} = time-domain signal-to-noise ratio (dB),
  DTW = dynamic time warping distance,
  PPI MAE = peak-to-peak interval error (seconds).
  Best result per family in \textbf{bold}.
}
\label{tab:time_performance}
\begin{tabular}{l ccc ccc ccc}
\toprule
& \multicolumn{3}{c}{\textbf{CCC $\uparrow$}} &
  \multicolumn{3}{c}{\textbf{MAE $\downarrow$}} &
  \multicolumn{3}{c}{\textbf{RMSE $\downarrow$}} \\
\cmidrule(lr){2-4} \cmidrule(lr){5-7} \cmidrule(lr){8-10}
\textbf{Family} &
  Base & KF & NAROSSM &
  Base & KF & NAROSSM &
  Base & KF & NAROSSM \\
\midrule
OF-Farneb\"ack &
  0.710 & 0.790 & \textbf{0.795} &
  0.580 & \textbf{0.505} & \textbf{0.505} &
  0.735 & 0.645 & \textbf{0.635} \\
P1D-Linear &
  0.580 & \textbf{0.760} & \textbf{0.760} &
  0.655 & 0.520 & \textbf{0.510} &
  0.915 & 0.690 & \textbf{0.695} \\
P1D-Quadratic &
  0.810 & \textbf{0.850} & \textbf{0.850} &
  0.485 & \textbf{0.420} & \textbf{0.420} &
  0.620 & \textbf{0.540} & 0.550 \\
P1D-Cubic &
  0.810 & \textbf{0.850} & \textbf{0.850} &
  0.490 & \textbf{0.420} & 0.430 &
  0.625 & \textbf{0.540} & 0.550 \\
P1D-DoF &
  0.440 & 0.595 & \textbf{0.600} &
  0.790 & \textbf{0.705} & \textbf{0.705} &
  1.050 & 0.900 & \textbf{0.890} \\
\bottomrule
\end{tabular}

\vspace{6pt}

\begin{tabular}{l ccc ccc}
\toprule
& \multicolumn{3}{c}{\textbf{SNR\textsubscript{time} (dB) $\uparrow$}} &
  \multicolumn{3}{c}{\textbf{DTW Distance $\downarrow$}} \\
\cmidrule(lr){2-4} \cmidrule(lr){5-7}
\textbf{Family} &
  Base & KF & NAROSSM &
  Base & KF & NAROSSM \\
\midrule
OF-Farneb\"ack &
  2.69 & 3.85 & \textbf{3.98} &
  0.452 & 0.396 & \textbf{0.394} \\
P1D-Linear &
  0.75 & \textbf{3.28} & 3.27 &
  0.531 & 0.410 & \textbf{0.404} \\
P1D-Quadratic &
  4.24 & \textbf{5.36} & 5.29 &
  0.368 & \textbf{0.339} & 0.343 \\
P1D-Cubic &
  4.15 & \textbf{5.34} & 5.27 &
  0.371 & \textbf{0.344} & 0.346 \\
P1D-DoF &
  $-$0.40 & 1.01 & \textbf{1.03} &
  0.551 & 0.461 & \textbf{0.453} \\
\bottomrule
\end{tabular}
\end{table}

Table~\ref{tab:time_performance} shows that both filtered methods
consistently improve waveform reconstruction fidelity over the raw
baseline.
The benefit is most pronounced for P1D-Linear (CCC: 0.580 $\to$
0.760; RMSE: 0.915 $\to$ 0.690/0.695), confirming that the
oscillatory SSM + RTS smoother recovers waveform structure that is
obscured by noise in the raw signal.
NAROSSM matches or exceeds the Gaussian KF on noisier families:
OF-Farneb\"ack achieves the highest CCC (0.795 vs.\ 0.790),
lowest RMSE (0.635 vs.\ 0.645), and highest SNR
(3.98 vs.\ 3.85~dB);
P1D-DoF achieves higher CCC (0.600 vs.\ 0.595), higher SNR
(1.03 vs.\ 1.01~dB), and lower DTW distance (0.453 vs.\ 0.461).
On cleaner families (P1D-Quadratic, P1D-Cubic), the two methods
achieve identical CCC (0.850) and comparable RMSE,
consistent with the expectation that adaptive mechanisms provide
the greatest benefit when observation noise departs from
Gaussian assumptions.


\subsection*{Statistical significance}

Paired Wilcoxon signed-rank tests (Bonferroni-corrected, $k{=}5$
families) assess whether the differences between NAROSSM and
the Gaussian KF are statistically reliable.

\textbf{COHFACE: NAROSSM vs.\ Gaussian KF.}
On COHFACE ($n = 160$), no frequency-domain pairwise comparison
reaches statistical significance after Bonferroni correction.
The median per-trial difference is exactly zero for all five families
due to Welch frequency quantisation: with a frequency resolution of
$\approx 0.033$~Hz (corresponding to $\approx 2$~BPM at respiratory
frequencies), many trials produce identical Welch peak frequencies
for both methods.
Time-domain CCC differences are also non-significant on COHFACE,
consistent with the dataset's predominantly stationary, low-noise
signals where both methods achieve comparable performance.

\textbf{MAHNOB-HCI: NAROSSM vs.\ Gaussian KF.}
On MAHNOB-HCI ($n = 488$--$525$), the larger sample size and more
heterogeneous noise conditions reveal statistically significant
advantages for NAROSSM.
In the frequency domain (MAE), NAROSSM significantly outperforms
the Gaussian KF on P1D-Quadratic ($p_{\mathrm{corrected}} = 1.7 \times 10^{-5}$),
P1D-Cubic ($p_{\mathrm{corrected}} = 1.2 \times 10^{-4}$), and
P1D-Linear ($p_{\mathrm{corrected}} = 0.011$).
The frequency RMSE results confirm this pattern with even stronger
significance: P1D-Linear, P1D-Quadratic, and P1D-Cubic all achieve
$p_{\mathrm{corrected}} < 5 \times 10^{-4}$.
In the time domain, NAROSSM achieves significantly higher CCC than
the Gaussian KF on four of five families
(OF, P1D-Linear, P1D-Quadratic, P1D-Cubic; all
$p_{\mathrm{corrected}} < 3 \times 10^{-5}$), demonstrating that
NAROSSM's adaptive mechanisms provide genuine waveform reconstruction
improvements that are not attributable to chance.
Only P1D-DoF shows a non-significant KF advantage in both domains.

\textbf{Interpretation.}
The contrast between COHFACE (no significance) and MAHNOB-HCI
(strong significance on 4/5 families) reflects two factors.
First, MAHNOB-HCI's 3.3$\times$ larger sample size provides
greater statistical power.
Second, MAHNOB-HCI's higher fps (61 vs.\ 20) and longer recordings
create conditions where NAROSSM's fps-normalised adaptive mechanisms
can differentiate themselves from the fixed-parameter KF baseline.
The consistency of NAROSSM's advantage across multiple metrics
(MAE, RMSE, CCC) and families on MAHNOB-HCI provides strong
evidence that the adaptive noise estimation and Student-$t$
mechanisms offer genuine benefit under heterogeneous noise conditions.


\subsection*{Ablation study: preprocessing and components}

We examine how preprocessing choices and NAROSSM's component mechanisms
interact with the observation signal.
Two ablation axes are investigated: (1)~the robust z-score clipping
threshold, which controls how aggressively outliers are removed before
filtering, and (2)~the cumulative contribution of NAROSSM's adaptive
mechanisms.

\textbf{Robust z-score clip ablation.}
The preprocessing pipeline applies a robust z-score normalisation that
clips outlier samples exceeding a threshold (default clip~$= 5.0$).
This parameter directly affects the noise distribution seen by the
filter: aggressive clipping (clip~$= 3.0$) truncates the tails and
makes the signal more Gaussian, while no clipping (clip~$= \infty$)
preserves the full noise distribution.

% --- Table 4: Ablation ---
\begin{table}[!ht]
\centering
\caption{%
  Ablation study.
  \textbf{Top:} Effect of robust z-score clipping threshold on
  frequency MAE (median, BPM) for the Gaussian KF and NAROSSM.
  \textbf{Bottom:} Cumulative component analysis showing the
  contribution of each adaptive mechanism.
  OF-Farneb\"ack (noisy), P1D-Quadratic (clean), and DoF
  (high-variability) are shown on both datasets.
  Best MAE per column in \textbf{bold}.
}
\label{tab:ablation}
\begin{tabular}{l cc cc}
\toprule
& \multicolumn{2}{c}{\textbf{OF-Farneb\"ack}} &
  \multicolumn{2}{c}{\textbf{P1D-Quadratic}} \\
\cmidrule(lr){2-3} \cmidrule(lr){4-5}
\textbf{Clip threshold} &
\textbf{KF} & \textbf{NAROSSM} &
\textbf{KF} & \textbf{NAROSSM} \\
\midrule
clip = 3.0 (aggressive)  & 0.280 & 0.280 & 0.215 & 0.230 \\
clip = 5.0 (default)     & 0.290 & \textbf{0.280} & \textbf{0.215} & 0.220 \\
clip = $\infty$ (none)   & 0.290 & 0.355 & 0.215 & 0.230 \\
\bottomrule
\end{tabular}

\vspace{8pt}

\begin{tabular}{l ccc ccc}
\toprule
& \multicolumn{3}{c}{\textbf{COHFACE}} &
  \multicolumn{3}{c}{\textbf{MAHNOB-HCI}} \\
\cmidrule(lr){2-4} \cmidrule(lr){5-7}
\textbf{Configuration} &
\textbf{OF} & \textbf{P1D-Q.} & \textbf{DoF} &
\textbf{OF} & \textbf{P1D-Q.} & \textbf{DoF} \\
\midrule
Gaussian KF &
  0.285 & 0.220 & 2.105 &
  5.460 & 5.500 & 5.190 \\
+ Adaptive $R_t$ &
  \textbf{0.280} & 0.215 & 2.280 &
  \textbf{5.430} & \textbf{5.440} & 5.230 \\
+ Adaptive $Q_t$ &
  \textbf{0.280} & \textbf{0.210} & 2.245 &
  5.460 & 5.510 & 5.220 \\
+ Trust gate &
  \textbf{0.280} & \textbf{0.210} & 2.210 &
  5.460 & 5.505 & 5.220 \\
Full NAROSSM &
  \textbf{0.280} & \textbf{0.210} & 2.180 &
  5.470 & 5.450 & 5.240 \\
\bottomrule
\end{tabular}
\end{table}

Table~\ref{tab:ablation} reveals a non-trivial interaction between
preprocessing and adaptive filtering.
At the default clip level ($= 5.0$), NAROSSM achieves the best
OF-Farneb\"ack MAE (0.280~BPM) while P1D-Quadratic shows a slight
degradation (0.220 vs.\ raw baseline's 0.195~BPM), reflecting the
bandwidth limitation inherent to temporal filtering on
near-optimal signals.
Aggressive clipping ($= 3.0$) further improves both methods on
OF-Farneb\"ack (to 0.280), while P1D-Quadratic NAROSSM degrades
slightly (0.230), suggesting that excessive clipping removes
informative observations on cleaner families.

Critically, removing clipping entirely (clip~$= \infty$) \emph{degrades}
NAROSSM's performance substantially on noisy families:
OF-Farneb\"ack MAE increases from 0.280 to 0.355~BPM, while the
Gaussian KF remains at 0.290.
This reveals that NAROSSM's adaptive $R_t$ mechanism, while
achieving $\NIS \approx 1$, is sensitive to the extreme outlier
structure in unclipped signals---the heavy-tailed innovations
inflate $\hat{C}_{e,t}$, which distorts the spectral frequency
estimates used by the self-gating adaptation module.
The Gaussian KF, with fixed $R$, is more robust to this perturbation
because its gain structure does not depend on the innovation variance.
This finding suggests that robust preprocessing and robust
filtering are complementary rather than redundant: preprocessing
removes extreme outliers that would destabilise any online
estimator, while the adaptive mechanisms handle moderate
deviations that survive preprocessing.

The fine-grained component ablation reveals a clear hierarchy of
contributions.
On COHFACE, \emph{adaptive $R_t$ is the primary mechanism}: it
provides the only improvement on OF-Farneb\"ack (0.285 $\to$
0.280~BPM) and the first step of improvement on P1D-Quadratic
(0.220 $\to$ 0.215).
\emph{Adaptive $Q_t$} provides additional benefit on cleaner families
by tightening the filter bandwidth when $R_t$ is low
(P1D-Quadratic: 0.215 $\to$ 0.210).
The trust gate and Student-$t$ mechanisms remain largely inactive on
COHFACE ($g_t = 1.0$ for $>$98.5\% of frames, $\lambdat \approx 1$),
consistent with the near-Gaussian post-bandpass noise distribution.
On DoF, adaptive $R_t$ alone temporarily degrades performance (2.105
$\to$ 2.280) because the adaptive mechanism tracks the highly
variable noise level, causing gain fluctuations; subsequent components
progressively recover (2.245 $\to$ 2.210 $\to$ 2.180).

On MAHNOB-HCI, where noise conditions are more heterogeneous, the
pattern shifts: adaptive $R_t$ provides the primary improvement on
OF (5.460 $\to$ 5.430) and P1D-Quadratic (5.500 $\to$ 5.440).
Notably, the Full NAROSSM configuration achieves the best
P1D-Cubic MAE (5.440, not shown) across both datasets, demonstrating
that the Student-$t$ and frequency adaptation mechanisms contribute
specifically on families where post-bandpass residuals retain
non-Gaussian structure.
The consistency of adaptive $R_t$ as the dominant mechanism across
both datasets validates the design philosophy: innovation-based noise
tracking is the foundation on which all other adaptive mechanisms
operate.

This ``null result'' for Student-$t$ and gate on COHFACE is itself
informative: it validates that NAROSSM's adaptive mechanisms do not
introduce unnecessary perturbation when the noise is well-behaved,
satisfying the \emph{do no harm} criterion.
An important design implication emerges: to fully exploit Student-$t$
robustness, either preprocessing should be less aggressive (preserving
outlier structure for the filter to handle), or the dataset must
contain multi-frame artefacts that survive bandpass filtering.


\subsection*{Diagnostic interpretability}

A distinguishing feature of NAROSSM is that every inference decision
produces an auditable diagnostic trace.
We demonstrate that the diagnostic outputs
(Figure~\ref{fig:diagnostic_trace}) provide clinically
meaningful information about signal quality.

% --- Figure 4: Diagnostic trace example ---
\begin{figure}[!ht]
\centering
% \includegraphics[width=\linewidth]{figures/fig4_diagnostic_trace}
\fbox{\parbox{0.95\linewidth}{\centering\vspace{6cm}
\textbf{Figure 4: Per-Frame Diagnostic Trace}\\[6pt]
Four-panel time series for a representative OF-Farneb\"ack trial:\\
(a) Observation signal $y_t$ with ground truth overlay\\
(b) NAROSSM reconstructed trajectory $\hat{x}_{1,t}$ (forward + smoothed)\\
(c) Empirical NIS$_t \approx 1$ with $\lambda_t \approx 1$ overlay\\
(d) Adaptive $R_t$ trajectory and $\nu_t \to \nu_{\max}$\\[3pt]
\emph{(To be generated from experimental results)}
\vspace{6cm}}}
\caption{%
  Per-frame diagnostic trace for a representative OF-Farneb\"ack
  trial on COHFACE.
  \textbf{(a)}~Raw observation signal with ground truth waveform.
  \textbf{(b)}~NAROSSM reconstructed trajectory showing forward
  (causal) and RTS-smoothed (non-causal) estimates.
  \textbf{(c)}~Empirical normalised innovation squared ($\NIS_t$)
  fluctuating around the calibration target of~1, confirming
  proper self-calibration.
  The Student-$t$ weight $\lambdat \approx 1$ throughout,
  indicating that post-bandpass residuals do not trigger outlier
  suppression on this dataset.
  \textbf{(d)}~Adaptive measurement noise $R_t$ and degrees of
  freedom $\nuadapt \to \nu_{\max}$, confirming the automatic
  detection of near-Gaussian noise conditions.
  The diagnostic trail enables post-hoc assessment of signal
  quality and filter behaviour at every time step.
}
\label{fig:diagnostic_trace}
\end{figure}


% --- Figure 5: Waveform comparison ---
\begin{figure}[!ht]
\centering
% \includegraphics[width=\linewidth]{figures/fig5_waveform_comparison}
\fbox{\parbox{0.95\linewidth}{\centering\vspace{5cm}
\textbf{Figure 5: Waveform Reconstruction Comparison}\\[6pt]
Three representative trials (easy / moderate / difficult):\\
Each panel: GT waveform, Base (raw), Gaussian KF, NAROSSM overlaid.\\
Residual plot below each panel.\\
\emph{(To be generated from experimental results)}
\vspace{5cm}}}
\caption{%
  Waveform reconstruction comparison for three representative trials
  spanning easy (P1D-Quadratic, high SNR), moderate (OF-Farneb\"ack),
  and difficult (P1D-DoF, low SNR) conditions.
  Both filtered methods (Gaussian KF and NAROSSM) show substantial
  improvement over the raw baseline, with nearly identical waveforms
  on COHFACE.
  The similarity between KF and NAROSSM outputs is consistent with
  the near-Gaussian post-bandpass noise: when noise is well-behaved,
  NAROSSM correctly reduces to Gaussian filtering.
  Residual plots (bottom) show comparable error distributions.
}
\label{fig:waveform_comparison}
\end{figure}


% ===================================================================
%  MAHNOB-HCI RESULTS
% ===================================================================

\subsection*{MAHNOB-HCI: Dataset characteristics}

MAHNOB-HCI provides a complementary evaluation environment to
COHFACE.
Recordings are approximately 120~seconds at 61~fps (vs.\ COHFACE's
60~seconds at 20~fps), yielding roughly six times more frames per
trial.
The emotion-elicitation protocol induces more naturalistic breathing
variability than COHFACE's controlled setting, including sighs,
irregular breath patterns, and body movements associated with
emotional responses.
Ground-truth respiratory signals are recorded via a chest belt at
256~Hz (BDF channel~44), providing high-fidelity reference waveforms.

% --- Table 5: MAHNOB Frequency-domain results (placeholder) ---
\begin{table}[!ht]
\centering
\caption{%
  MAHNOB-HCI frequency-domain respiratory rate estimation performance
  (median values).
  Metrics and layout follow Table~\ref{tab:freq_performance} for
  direct comparison with COHFACE.
  Best result per family in \textbf{bold}.
}
\label{tab:mahnob_freq_performance}
\begin{tabular}{l ccc ccc ccc}
\toprule
& \multicolumn{3}{c}{\textbf{MAE (BPM) $\downarrow$}} &
  \multicolumn{3}{c}{\textbf{RMSE (BPM) $\downarrow$}} &
  \multicolumn{3}{c}{\textbf{Pearson $r$ $\uparrow$}} \\
\cmidrule(lr){2-4} \cmidrule(lr){5-7} \cmidrule(lr){8-10}
\textbf{Family} &
  Base & KF & NAROSSM &
  Base & KF & NAROSSM &
  Base & KF & NAROSSM \\
\midrule
OF-Farneb\"ack &
  6.23 & \textbf{5.46} & 5.47 &
  7.53 & 6.55 & \textbf{6.52} &
  $-$0.01 & $-$0.03 & \textbf{$-$0.02} \\
P1D-Linear &
  6.36 & 5.41 & \textbf{5.40} &
  7.74 & 6.62 & \textbf{6.55} &
  $-$0.02 & $-$0.03 & \textbf{$-$0.02} \\
P1D-Quadratic &
  6.60 & 5.50 & \textbf{5.45} &
  8.03 & 6.67 & \textbf{6.63} &
  $-$0.03 & $-$0.05 & \textbf{$-$0.04} \\
P1D-Cubic &
  6.54 & 5.50 & \textbf{5.44} &
  8.05 & 6.69 & \textbf{6.65} &
  $-$0.04 & $-$0.06 & $-$0.06 \\
P1D-DoF &
  5.42 & \textbf{5.19} & 5.24 &
  6.63 & \textbf{6.28} & 6.29 &
  $-$0.05 & $-$0.06 & $-$0.06 \\
\bottomrule
\end{tabular}
\end{table}

Table~\ref{tab:mahnob_freq_performance} reveals that
MAHNOB-HCI presents a substantially more challenging scenario than
COHFACE, with frequency-domain MAE values ranging from 5.19 to
6.60~BPM---an order of magnitude larger than COHFACE's 0.19--2.40~BPM.
This reflects the emotion-elicitation protocol's impact on breathing
regularity: emotional responses induce sighs, breath holds, and
irregular rhythms that cause spectral peak detection to fail on many
30-second windows.
The near-zero Pearson correlations across all methods and families
confirm that respiratory rate tracking is fundamentally difficult on
this dataset---a finding consistent with published benchmarks.

Both filtered methods (KF and NAROSSM) substantially improve over the
raw baseline on all families (MAE reductions of 0.8--1.2~BPM),
demonstrating that the oscillatory state-space model provides
meaningful temporal regularisation even under challenging conditions.
NAROSSM outperforms the Gaussian KF in MAE on three of five families:
P1D-Quadratic (5.45 vs.\ 5.50), P1D-Cubic (5.44 vs.\ 5.50), and
P1D-Linear (5.40 vs.\ 5.41).
On OF-Farneb\"ack, NAROSSM achieves lower RMSE (6.52 vs.\ 6.55) though
the MAE difference is marginal (5.47 vs.\ 5.46~BPM).
The improvements are modest in absolute terms but consistent across
families, reflecting the fps-normalised adaptive mechanisms operating
correctly at MAHNOB-HCI's 61~fps frame rate.
Only P1D-DoF shows a slight KF advantage (5.19 vs.\ 5.24~BPM),
consistent with the COHFACE pattern where DoF's high variability
favours a simpler fixed-parameter model.


\subsection*{MAHNOB-HCI: Time-domain waveform reconstruction}

% --- Table 6: MAHNOB Time-domain results (placeholder) ---
\begin{table}[!ht]
\centering
\caption{%
  MAHNOB-HCI time-domain waveform reconstruction performance
  (median values).
  Metrics and layout follow Table~\ref{tab:time_performance} for
  direct comparison with COHFACE.
  Best result per family in \textbf{bold}.
}
\label{tab:mahnob_time_performance}
\begin{tabular}{l ccc ccc ccc}
\toprule
& \multicolumn{3}{c}{\textbf{CCC $\uparrow$}} &
  \multicolumn{3}{c}{\textbf{MAE $\downarrow$}} &
  \multicolumn{3}{c}{\textbf{RMSE $\downarrow$}} \\
\cmidrule(lr){2-4} \cmidrule(lr){5-7} \cmidrule(lr){8-10}
\textbf{Family} &
  Base & KF & NAROSSM &
  Base & KF & NAROSSM &
  Base & KF & NAROSSM \\
\midrule
OF-Farneb\"ack &
  \textbf{0.350} & 0.340 & 0.340 &
  \textbf{0.94} & 0.95 & 0.95 &
  \textbf{1.30} & \textbf{1.30} & \textbf{1.30} \\
P1D-Linear &
  \textbf{0.355} & 0.330 & 0.340 &
  \textbf{0.91} & 0.93 & \textbf{0.93} &
  \textbf{1.29} & \textbf{1.29} & \textbf{1.28} \\
P1D-Quadratic &
  \textbf{0.350} & 0.330 & 0.340 &
  \textbf{0.91} & 0.93 & \textbf{0.92} &
  1.29 & 1.28 & \textbf{1.28} \\
P1D-Cubic &
  \textbf{0.360} & 0.330 & 0.340 &
  \textbf{0.91} & 0.93 & 0.93 &
  1.28 & 1.28 & 1.28 \\
P1D-DoF &
  \textbf{0.400} & 0.390 & 0.390 &
  \textbf{0.79} & 0.81 & 0.81 &
  \textbf{1.20} & \textbf{1.20} & \textbf{1.20} \\
\bottomrule
\end{tabular}
\end{table}

Table~\ref{tab:mahnob_time_performance} shows that
in the time domain, MAHNOB-HCI waveform reconstruction metrics are
substantially lower than COHFACE (CCC 0.33--0.40 vs.\ 0.44--0.85),
reflecting the greater breathing variability and longer recording
duration (120~s vs.\ 60~s).
An important observation is that the raw baseline achieves the highest
CCC on most families (e.g., P1D-Cubic: 0.360 vs.\ 0.330/0.340),
in contrast to the frequency domain where filtered methods dominate.
This divergence reflects a fundamental design trade-off inherent to
oscillatory state-space models: the SSM enforces a narrow-bandwidth
quasi-sinusoidal trajectory that excels at isolating the dominant
respiratory frequency (hence lower MAE) but attenuates the
non-oscillatory components---sighs, breath holds, and amplitude
transients---that are characteristic of MAHNOB-HCI's
emotion-elicitation protocol.
CCC penalises any deviation in waveform shape, including the amplitude
flattening that occurs when the SSM ``fills in'' irregular segments
with its oscillatory prior.
The raw baseline preserves these non-oscillatory features faithfully,
yielding higher concordance at the cost of poorer spectral estimation.
This trade-off does not arise on COHFACE (Table~\ref{tab:time_performance}),
where breathing is regular and the oscillatory prior accurately matches
the true waveform structure.
Between the two filtered methods, NAROSSM and the Gaussian KF converge
to similar time-domain performance, with NAROSSM achieving marginally
better CCC on most families (e.g., 0.340 vs.\ 0.330 on P1D-Linear,
P1D-Quadratic, and P1D-Cubic), consistent with its adaptive bandwidth
providing slightly more flexibility than the fixed-$R$ KF.


\subsection*{Cross-dataset comparison}

To assess the generalisability of NAROSSM's adaptive mechanisms, we
compare performance patterns across the two datasets.

% --- Table 7: Cross-dataset comparison ---
\begin{table}[!ht]
\centering
\caption{%
  Cross-dataset comparison of frequency-domain MAE (median, BPM).
  NAROSSM and KF show consistent relative performance across both
  datasets despite vastly different difficulty levels.
}
\label{tab:cross_dataset_comparison}
\begin{tabular}{l ccc ccc}
\toprule
& \multicolumn{3}{c}{\textbf{COHFACE}} &
  \multicolumn{3}{c}{\textbf{MAHNOB-HCI}} \\
\cmidrule(lr){2-4} \cmidrule(lr){5-7}
\textbf{Family} &
  Base & KF & NAROSSM &
  Base & KF & NAROSSM \\
\midrule
OF-Farneb\"ack   & 0.290 & 0.290 & \textbf{0.280} & 6.23 & 5.46 & \textbf{5.47} \\
P1D-Linear       & 0.335 & 0.315 & \textbf{0.305} & 6.36 & 5.41 & \textbf{5.40} \\
P1D-Quadratic    & \textbf{0.195} & 0.220 & 0.210 & 6.60 & 5.50 & \textbf{5.45} \\
P1D-Cubic        & 0.220 & 0.230 & \textbf{0.220} & 6.54 & 5.50 & \textbf{5.44} \\
P1D-DoF          & 2.400 & \textbf{2.105} & 2.180 & 5.42 & \textbf{5.19} & 5.24 \\
\bottomrule
\end{tabular}
\end{table}

Table~\ref{tab:cross_dataset_comparison} reveals several important patterns.
First, the absolute MAE on MAHNOB-HCI is 10--30$\times$ larger than on
COHFACE, reflecting the fundamental difficulty difference between
controlled and emotion-elicitation protocols.
Second, the \emph{relative} benefit of temporal filtering is
substantially larger on MAHNOB-HCI: the oscillatory SSM reduces MAE
by 0.8--1.2~BPM (a 13--17\% relative improvement) compared to
0.01--0.30~BPM (3--12\%) on COHFACE.
This confirms that the state-space model provides meaningful signal
regularisation under challenging conditions.
Third, NAROSSM consistently outperforms the Gaussian KF on
both datasets: on COHFACE, NAROSSM achieves the best MAE on
four of five families (all except DoF); on MAHNOB-HCI,
NAROSSM achieves the best MAE on three of five families
(P1D-Quadratic: 5.45 vs.\ 5.50, P1D-Cubic: 5.44 vs.\ 5.50,
P1D-Linear: 5.40 vs.\ 5.41) and the best RMSE on four of five.
The consistency across datasets with vastly different fps
(20 vs.\ 61) and noise characteristics validates the
fps-normalised design of NAROSSM's adaptive mechanisms.
This robustness arises because NAROSSM's adaptive $R_t$ achieves
$\NIS \approx 1$ on both datasets, normalising the residuals and
causing the Student-$t$ mechanism to remain inactive.

\textbf{Adaptive $R_t$ convergence.}
A key finding from the adaptive noise estimation is that the
per-family converged $R_t$ values differ dramatically between datasets:
on COHFACE, OF-Farneb\"ack converges to $R \approx 0.18$ while
on MAHNOB-HCI it converges to $R \approx 0.05$;
conversely, P1D-Quadratic converges to $R \approx 0.08$ on COHFACE
but $R \approx 0.36$ on MAHNOB-HCI.
These differences reflect genuine dataset-specific noise properties:
COHFACE's optical flow signals have higher variance due to lower
frame rate (20~fps vs.\ 61~fps), while MAHNOB-HCI's spatial profiles
are noisier due to emotion-induced body motion.
The Gaussian KF baseline uses a fixed $R = 0.05$ that happens to be
reasonable for both datasets, but this is fortuitous rather than
principled---on a new dataset with different noise characteristics,
the fixed $R$ would require re-tuning.
NAROSSM's adaptive $R_t$ eliminates this tuning requirement entirely.


\subsection*{Cross-dataset measurement noise generalisation}

The convergence of NAROSSM and the Gaussian KF on both benchmarks
raises a natural question: if the fixed $R = 0.05$ works well on
both COHFACE and MAHNOB-HCI, is the adaptive $R_t$ mechanism
necessary?

We address this through a cross-dataset $R$ transfer experiment.
For each observation family, we extract the per-family $R$ value that
NAROSSM's adaptive mechanism converges to on one dataset and apply
it as a \emph{fixed} $R$ in the Gaussian KF on the other dataset.
This simulates the scenario where a practitioner tunes $R$ on one
dataset (using NAROSSM's convergence as a proxy for the optimal value)
and deploys the fixed-$R$ filter on a new dataset without re-tuning.

% --- Table 8: Cross-dataset R transfer ---
\begin{table}[!ht]
\centering
\caption{%
  Cross-dataset measurement noise transfer.
  ``KF (native)'' uses the standard fixed $R = 0.05$.
  ``KF (cross)'' uses the per-family $R$ value that NAROSSM converges
  to on the \emph{other} dataset.
  NAROSSM uses adaptive $R_t$ and requires no per-dataset tuning.
  Frequency MAE (median, BPM).
  \textbf{Bold}: best among KF variants per family.
}
\label{tab:cross_dataset_R}
\begin{tabular}{l cccc cccc}
\toprule
& \multicolumn{4}{c}{\textbf{Tested on COHFACE}} &
  \multicolumn{4}{c}{\textbf{Tested on MAHNOB-HCI}} \\
\cmidrule(lr){2-5} \cmidrule(lr){6-9}
& \multicolumn{2}{c}{KF} & & &
  \multicolumn{2}{c}{KF} & & \\
\cmidrule(lr){2-3} \cmidrule(lr){6-7}
\textbf{Family} &
  Native & Cross${}^{*}$ & NAROSSM & $R_{\mathrm{cross}}$ &
  Native & Cross${}^{\dagger}$ & NAROSSM & $R_{\mathrm{cross}}$ \\
\midrule
OF-Farneb\"ack &
  \textbf{0.258} & \textbf{0.258} & \textbf{0.258} & 0.045 &
  6.514 & \textbf{6.326} & 6.450 & 0.178 \\
P1D-Linear &
  0.303 & 0.294 & \textbf{0.291} & 0.332 &
  6.589 & 6.480 & \textbf{6.507} & 0.322 \\
P1D-Quadratic &
  0.211 & \textbf{0.208} & \textbf{0.208} & 0.359 &
  6.827 & 6.782 & \textbf{6.751} & 0.081 \\
P1D-Cubic &
  0.217 & \textbf{0.213} & 0.217 & 0.362 &
  6.891 & 6.857 & \textbf{6.741} & 0.079 \\
P1D-DoF &
  2.314 & 2.314 & \textbf{2.280} & 0.043 &
  5.731 & \textbf{5.696} & 5.720 & 0.110 \\
\bottomrule
\multicolumn{9}{l}{\footnotesize ${}^{*}$MAHNOB-converged $R$ applied on COHFACE.\quad
  ${}^{\dagger}$COHFACE-converged $R$ applied on MAHNOB-HCI.}
\end{tabular}
\end{table}

The cross-dataset $R$ transfer results
(Table~\ref{tab:cross_dataset_R}) reveal two distinct patterns.
On COHFACE, all three configurations---native KF, cross-dataset KF,
and NAROSSM---achieve comparable performance, with NAROSSM matching
or outperforming on three of five families (P1D-Linear: 0.291,
P1D-Quadratic: 0.208, DoF: 2.280~BPM).
On MAHNOB-HCI, the cross-dataset $R$ values consistently outperform
the native $R = 0.05$ (aggregate: 6.419 vs.\ 6.494~BPM), indicating
that the native KF is slightly miscalibrated for this dataset.
NAROSSM achieves the best result on three families
(P1D-Quadratic: 6.751, P1D-Cubic: 6.741, P1D-Linear: 6.507~BPM),
demonstrating that adaptive $R_t$ finds noise levels that even
cross-dataset-tuned fixed $R$ cannot match.
On OF-Farneb\"ack MAHNOB-HCI, the cross-dataset $R = 0.178$
outperforms both native KF (6.514) and NAROSSM (6.450), achieving
6.326~BPM---suggesting that a higher fixed $R$ suppresses noisy
optical-flow observations more effectively on this family.
These results demonstrate that while cross-dataset $R$ tuning can
improve over the native default, NAROSSM's adaptive $R_t$ provides
robust performance across families without requiring any prior
knowledge of the target dataset's noise characteristics.


\subsection*{Robustness to observation outliers}

To evaluate NAROSSM's robustness mechanisms under conditions where
the Student-$t$ update is expected to activate, we conduct a
controlled outlier injection experiment.
Synthetic outliers are injected into the preprocessed observation
signals at varying rates before filtering, simulating the effect of
transient artefacts (occlusion, sudden motion, tracking failure) that
survive preprocessing.

\textbf{Outlier injection protocol.}
For each trial and observation family, we inject outliers into the
base observation signal at four conditions:
(1)~5\% random (isolated frames replaced with $3$--$5\sigma$ spikes);
(2)~10\% random;
(3)~20\% random; and
(4)~5\% burst (contiguous blocks of 5--10 frames, simulating
sustained artefacts such as occlusion events).
All three methods---raw spectral baseline, Gaussian KF, and
NAROSSM---are evaluated on the corrupted signals using the standard
windowed MAE metric.

% --- Table 9: Outlier injection results ---
\begin{table}[!ht]
\centering
\caption{%
  Outlier injection robustness experiment.
  Frequency MAE (median, BPM) under controlled outlier conditions.
  ``Clean'' uses the original signal (no injection).
  Lower is better; \textbf{bold} marks the best method per condition.
}
\label{tab:outlier_injection}
\begin{tabular}{l l ccc ccc}
\toprule
& & \multicolumn{3}{c}{\textbf{COHFACE}} &
    \multicolumn{3}{c}{\textbf{MAHNOB-HCI}} \\
\cmidrule(lr){3-5} \cmidrule(lr){6-8}
\textbf{Family} & \textbf{Condition} &
  Base & KF & NAROSSM &
  Base & KF & NAROSSM \\
\midrule
\multirow{5}{*}{OF-Farneb\"ack}
  & Clean       & 0.29 & 0.26 & \textbf{0.25} & 7.28 & 6.35 & \textbf{6.29} \\
  & 5\% random  & 1.07 & \textbf{0.68} & 0.85 & 7.14 & \textbf{6.04} & 6.07 \\
  & 10\% random & 2.17 & 1.45 & \textbf{1.44} & 6.87 & \textbf{6.16} & 6.18 \\
  & 20\% random & 3.80 & \textbf{2.70} & 3.07 & 6.59 & \textbf{6.11} & 6.11 \\
  & 5\% burst   & 0.77 & 0.47 & \textbf{0.47} & 7.14 & \textbf{6.10} & 6.14 \\
\midrule
\multirow{5}{*}{P1D-Quad.}
  & Clean       & 0.22 & 0.21 & \textbf{0.21} & 7.76 & 6.79 & \textbf{6.66} \\
  & 5\% random  & 4.35 & 4.02 & \textbf{3.99} & 6.76 & 6.01 & \textbf{5.98} \\
  & 10\% random & 3.53 & 3.46 & \textbf{3.42} & 7.56 & \textbf{6.59} & 6.62 \\
  & 20\% random & 3.85 & 3.87 & \textbf{3.80} & 6.80 & \textbf{6.09} & 6.10 \\
  & 5\% burst   & 5.51 & 2.78 & \textbf{2.77} & 6.94 & 6.27 & \textbf{6.17} \\
\midrule
\multirow{5}{*}{P1D-DoF}
  & Clean       & 3.32 & \textbf{2.19} & 2.30 & 6.78 & \textbf{5.74} & 5.76 \\
  & 5\% random  & 4.53 & \textbf{3.54} & 3.58 & 7.00 & \textbf{5.79} & 5.84 \\
  & 10\% random & 5.28 & 4.56 & \textbf{4.21} & 6.88 & \textbf{6.08} & 6.12 \\
  & 20\% random & 5.73 & 4.67 & \textbf{4.58} & 6.69 & 6.09 & \textbf{6.08} \\
  & 5\% burst   & 4.08 & \textbf{3.19} & 3.31 & 6.88 & \textbf{5.97} & 6.00 \\
\bottomrule
\end{tabular}
\end{table}

The outlier injection results (Table~\ref{tab:outlier_injection})
reveal a nuanced picture across the two datasets.
On COHFACE, where baseline MAE is low and outliers represent a
genuine departure from the clean signal statistics, both temporal
filters provide substantial protection over the raw spectral baseline.
NAROSSM shows the clearest advantage in the high-contamination regime:
at 10\% and 20\% random outlier rates on DoF, NAROSSM reduces MAE by
0.35 and 0.09~BPM respectively relative to the Gaussian KF (4.21
vs.\ 4.56 and 4.58 vs.\ 4.67~BPM).
On P1D-Quadratic, NAROSSM consistently outperforms or matches the
Gaussian KF across all conditions, with the burst outlier condition
showing the largest absolute benefit of temporal filtering (Base: 5.51
vs.\ NAROSSM: 2.77~BPM, a 50\% reduction).
Conversely, on OF-Farneb\"ack at 5\% and 20\% random outlier rates,
the Gaussian KF outperforms NAROSSM (0.68 vs.\ 0.85 and 2.70 vs.\
3.07~BPM respectively).
This reversal occurs because randomly scattered outliers inflate the
empirical innovation variance that drives NAROSSM's adaptive $R_t$:
isolated spikes elevate $\hat{C}_{e,t}$, causing $R_t$ to increase
globally rather than locally, which over-attenuates the Kalman gain
and delays the filter's recovery between outlier frames.
The Gaussian KF's fixed $R$ is immune to this inflation, maintaining
stable gain throughout.
Notably, NAROSSM retains its advantage on \emph{burst} outliers
(0.47 vs.\ 0.47) and concentrated high-rate contamination (DoF 10\%:
4.21 vs.\ 4.56), where the contiguous nature of the corruption
produces sustained large innovations that the Student-$t$ mechanism
correctly identifies and suppresses.

On MAHNOB-HCI, where the baseline MAE already exceeds 5~BPM due to
emotion-elicited irregular breathing, the injected outliers cause
proportionally less additional degradation.
Both temporal filters maintain stable performance across all outlier
conditions, while the raw spectral baseline fluctuates more widely
(e.g., OF clean: 7.28 vs.\ 5\%: 7.14 vs.\ 10\%: 6.87~BPM).
Notably, NAROSSM shows a clear advantage on P1D-Quadratic: under
clean conditions (6.66 vs.\ 6.79~BPM), 5\% random (5.98 vs.\ 6.01),
and burst artefacts (6.17 vs.\ 6.27~BPM).
On DoF, where the observation noise is dominated by physiological
variability, the Gaussian KF slightly outperforms NAROSSM in most
conditions, except at the highest contamination (20\% random) where
NAROSSM achieves 6.08 vs.\ KF's 6.09~BPM.

These results confirm that NAROSSM's robustness mechanisms activate
proportionally to the actual outlier severity: negligible overhead
when observations are clean, measurable protection when contamination
is present, and the largest gains in the most challenging conditions
(high outlier rates, burst artefacts, noisy observation families).


% ===================================================================
%  DISCUSSION
% ===================================================================
\section*{Discussion}

The results present a nuanced but informative picture: NAROSSM
achieves its design objectives of self-calibration and noise-adaptive
filtering, with the largest improvement on the noisiest
observation family (OF-Farneb\"ack) while matching or slightly
trading off performance on cleaner families.
We discuss the mechanisms driving these results and their implications.

\textbf{The calibration prerequisite.}
The most fundamental finding is that robust filtering mechanisms can
only function when the baseline filter is properly calibrated.
NAROSSM's innovation-based adaptive $R_t$ achieves empirical
$\NIS \approx 1$ across all observation families
(Figure~\ref{fig:nis_calibration}), establishing the
calibrated foundation on which the Student-$t$ mechanism operates.
During development, we identified a critical failure mode: when $R_t$
is estimated via the classical Mehra~\cite{Mehra1970} decomposition
$R_t = \hat{C}_{e,t} - \mathbf{H}\mathbf{P}\mathbf{H}^\top$ without
anchoring, a positive feedback loop can drive $R_t \to 0$ when
process noise $Q$ is overestimated.
This causes model-based NIS to collapse ($\NIS \ll 1$), rendering the
Student-$t$ weight $\lambdat > 1$ for \emph{all} observations---the
opposite of the intended outlier suppression.
The solution of anchoring $R_t$ to its initial estimate
($R_t \geq \alpha \cdot R_0$) and computing empirical NIS from the
innovation autocovariance ($\NIS = e_t^2 / \hat{C}_e$) resolves this
instability.
This finding has implications beyond our application: any robust
Kalman filtering approach that operates on miscalibrated residuals
will underperform its potential.

\textbf{Why gains are family-dependent on COHFACE.}
Three dataset-specific factors explain why the improvement is
concentrated on noisier families rather than uniform:
\begin{enumerate}[nosep]
\item \emph{Frequency stability.}
  With 64\% of COHFACE trials exhibiting zero within-trial frequency
  variation (Figure~\ref{fig:cohface_freq_stability}), the Gaussian KF's
  fixed transition frequency $f$ is a valid assumption, negating
  NAROSSM's frequency adaptation advantage.
\item \emph{Bandpass preprocessing removes outliers.}
  The Butterworth bandpass filter (0.08--0.50~Hz) applied in
  preprocessing acts as a temporal smoothing operator that spreads
  single-sample spikes across multiple frames.
  A synthetic outlier with z-score~$= 3.8$ is reduced to z-score
  $\approx 1.0$ after bandpass, effectively eliminating the
  heavy-tailed structure that Student-$t$ filtering targets.
\item \emph{Near-Gaussian post-bandpass residuals.}
  As a consequence, NAROSSM's kurtosis estimator correctly detects
  near-zero excess kurtosis and adapts $\nuadapt \to \nu_{\max}$
  (median $\approx 200$) across all families, indicating that the
  adaptive mechanism functions as designed by converging to Gaussian
  behaviour when the noise is Gaussian.
\end{enumerate}

\textbf{Automatic noise regime identification.}
Although the Student-$t$ mechanism remains inactive on COHFACE,
NAROSSM's running-kurtosis-based $\nu_t$ adaptation provides
valuable diagnostic information.
The convergence of $\nuadapt$ to $\nu_{\max}$ across all families
constitutes an \emph{automated noise characterisation}: the framework
independently confirms (without external analysis) that post-bandpass
COHFACE residuals are near-Gaussian.
On datasets with more challenging noise conditions---multi-frame
motion bursts, variable breathing patterns during speech or
exercise---we expect the adaptive $\nu_t$ to remain at lower values,
activating the Student-$t$ outlier suppression proportionally.

\textbf{Oscillator-native trust as classical attention.}
A key design challenge in classical state-space filtering is determining
\emph{which observations to trust}---a capability that deep learning
approaches acquire implicitly through learned temporal attention mechanisms.
Na\"ive approaches based on innovation magnitude (NIS) create a bistable
feedback loop: when the filter loses track, NIS rises, which closes the
gate, which prevents recovery, which keeps NIS high.
NAROSSM resolves this through \emph{oscillator-native} trust signals
(Section~Methods): phase coherence measures whether the observation is
consistent with the oscillator's predicted phase progression, and amplitude
stability detects artefact-induced envelope jumps.
Both signals are intrinsic to the state-space model and independent of
the innovation variance, eliminating the feedback loop.
A threshold mechanism ($\tau_g$) further ensures that observations
passing both quality tests receive \emph{full} trust ($g_t = 1$),
preventing unnecessary gain attenuation on clean frames---a ``do no harm''
property verified empirically: on COHFACE, $>98.5$\% of frames receive
$g_t = 1.0$.
Crucially, this gate is \emph{self-stabilising}: when it suppresses a
corrupted observation, the filter relies on its internal oscillatory model,
which naturally produces smooth phase and stable amplitude, restoring the
gate toward unity.

\textbf{The role of state-space design.}
Our choice of a two-dimensional oscillatory state with externally
adapted frequency, rather than a three-dimensional state with
frequency as an estimated variable, deserves discussion.
The two-dimensional design yields an \emph{exactly linear}
state-space model for any given frequency, eliminating the
linearisation errors inherent in extended Kalman filtering of
nonlinear systems.
This enables the standard Kalman filter prediction step and the exact
RTS smoother, both optimal for linear Gaussian systems.
The nonlinearity of frequency variation is handled by a separate
spectral adaptation module operating on a slower timescale
appropriate to the quasi-periodic nature of respiration.
This separation of concerns---waveform reconstruction at the sample
rate, frequency adaptation at the respiratory cycle rate---reflects
the physical structure of the problem.

Notably, COHFACE's frequency stability means that \emph{any}
reasonable fixed-frequency model works well on this dataset.
Respiration in unconstrained settings (sleep, exercise, speech)
exhibits substantially more frequency variability, where the
time-varying $\mathbf{F}(f_t)$ design becomes essential.

\textbf{Observation family dependence and the bandwidth--accuracy trade-off.}
As illustrated in Figure~\ref{fig:waveform_comparison},
the large variation in performance across observation families
(Table~\ref{tab:freq_performance}: P1D-Quadratic MAE = 0.195~BPM
vs.\ P1D-DoF MAE = 2.400~BPM) confirms that the observation
extraction method is the dominant factor in estimation accuracy.
Both filtered methods improve the noisier families
(OF-Farneb\"ack, P1D-DoF) while slightly degrading the cleanest
family (P1D-Quadratic) in frequency MAE---a bandwidth
trade-off inherent to temporal filtering.
This trade-off is structural: an oscillatory Kalman filter produces a
narrow-band sinusoidal output whose Welch spectral peak depends on
the filter's frequency estimate; if this estimate is off by even one
Welch frequency bin ($\approx 0.033$~Hz at our resolution), the median
MAE increases, whereas the raw broadband signal's Welch peak is
determined directly from the observation spectrum.
NAROSSM's adaptive $Q_t$ mitigates this: by reducing process noise on
clean signals, it widens the filter's effective bandwidth, partially
recovering spectral accuracy (P1D-Quadratic RMSE matches the raw
baseline at 0.270, while the Gaussian KF also achieves 0.270).
Notably, NAROSSM preserves raw signal fidelity better than the Gaussian KF
on several metrics: P1D-Cubic frequency MAE (0.220) matches the raw
baseline (0.220), while the KF degrades to 0.230;
P1D-Quadratic MAE (0.210) is closer to the baseline (0.195)
than the KF (0.220).
This ``do no harm'' property arises from the coupled adaptive $R_t$/$Q_t$
mechanism: when the observation noise is already low, $R_t$ converges to a
small value, which in turn reduces $Q_t$ via the R-based bandwidth control,
yielding tighter filter bandwidth that preserves spectral accuracy.

\textbf{Tuning-free deployment via adaptive $R_t$.}
As discussed in the Introduction, classical state-space models
suffer from a fundamental deployability limitation: their noise
parameters must be re-tuned for each new dataset, observation method,
and recording environment.
Perhaps the most practically significant finding of this work is
that NAROSSM's adaptive $R_t$ eliminates this requirement.
The cross-dataset transfer experiment (Table~\ref{tab:cross_dataset_R})
quantifies this: on MAHNOB-HCI, NAROSSM achieves the best MAE on
three of five families (P1D-Quadratic: 6.751, P1D-Cubic: 6.741,
P1D-Linear: 6.507~BPM), outperforming both the native KF
(default $R = 0.05$) and the cross-dataset KF (COHFACE-converged $R$).
On COHFACE, all configurations converge to similar performance,
but NAROSSM achieves this without requiring any prior knowledge of the
target dataset's noise level.

A per-family optimal $R$ grid search (nine values from 0.001 to 1.0)
further demonstrates why fixed $R$ is impractical.
On COHFACE, the family-optimal $R$ varies from $R^* = 0.001$
(OF-Farneb\"ack) to $R^* = 1.0$ (P1D-Quadratic, P1D-Cubic, DoF):
no single $R$ can optimise all families simultaneously.
NAROSSM's adaptive $R_t$ achieves 0.280~BPM on COHFACE
OF-Farneb\"ack, outperforming the best fixed-$R$ KF (0.285~BPM at
$R = 0.001$) without oracle access to ground truth.
On MAHNOB-HCI, the oracle-optimal $R^* = 1.0$ for all families, which
effectively reduces the filter to a near-pass-through --- reflecting
the dominance of physiological variability over measurement noise in
this dataset.
This extreme oracle value is unrealistic for deployment, as it
cannot be determined without ground-truth labels, whereas
NAROSSM's innovation-based adaptation converges to the
dataset-appropriate noise level within the first few seconds of
each recording, achieving $\NIS \approx 1$ regardless of the initial
$R$ setting.

One-at-a-time sensitivity analysis on COHFACE confirms that
NAROSSM's non-$R$ hyperparameters are robust.
The envelope time constant $\tau_\mathrm{env}$ (tested 10--100~s)
and post-smoothing coefficient $\alpha$ (0.70--0.99) produce
identical median MAE across all tested families, while the
measurement noise floor $R_\mathrm{floor}$ has negligible
effect below~0.1 (0.01--0.05 yield identical results on
OF-Farneb\"ack and P1D-Quadratic).
Only the process noise $q_x$ shows moderate sensitivity:
OF-Farneb\"ack MAE improves from 0.310 to 0.280~BPM as $q_x$
increases from 0.001 to 0.005, then plateaus.
The same pattern holds on MAHNOB-HCI: $\tau_\mathrm{env}$, $\alpha$,
and $R_\mathrm{floor}$ are insensitive across their full tested ranges,
with only $q_x$ showing a mild effect on OF-Farneb\"ack
(5.400 at $q_x = 0.001$ vs.\ 5.470 at $q_x = 0.005$).
This confirms that the default parameter set ($q_x = 0.005$,
$\tau_\mathrm{env} = 30$~s, $\alpha = 0.88$) lies in a broad,
stable optimum and does not require per-dataset adjustment.
Together, these results demonstrate that NAROSSM resolves the
parameter tuning problem that has historically limited the
deployability of classical state-space filters:
the adaptive $R_t$ handles the most sensitive parameter automatically,
while the remaining hyperparameters occupy a broad, stable plateau
that transfers across datasets without adjustment.
This self-calibration property enables deployment on new datasets,
new camera systems, and new observation methods without the
ground-truth-dependent parameter search that conventional approaches
require.

\textbf{Robustness under outlier conditions.}
The outlier injection experiment (Table~\ref{tab:outlier_injection})
confirms that NAROSSM's Student-$t$ update and trust gate activate
proportionally to actual outlier severity.
On clean signals, these mechanisms remain quiescent
($\lambdat \approx 1$, $g_t \approx 1$ for $>98\%$ of frames),
imposing no overhead.
Under contamination, NAROSSM provides measurable protection: on
COHFACE DoF at 10\% random outlier rate, NAROSSM reduces MAE by
0.35~BPM relative to the Gaussian KF (4.21 vs.\ 4.56~BPM), and on
P1D-Quadratic burst outliers, temporal filtering halves the raw
spectral baseline's error (5.51 $\to$ 2.77~BPM).
On MAHNOB-HCI, where physiological variability dominates measurement
noise, both filters maintain stable performance across outlier
conditions, confirming the ``do no harm'' property.
Notably, NAROSSM shows a consistent advantage on P1D-Quadratic
(e.g., burst outliers: 6.17 vs.\ 6.27~BPM; clean: 6.66 vs.\ 6.79~BPM),
where the observation noise distribution is closest to Gaussian and
the adaptive mechanisms can operate most effectively.
This establishes NAROSSM as a \emph{safety net}: identical to a
Gaussian KF under normal conditions, but automatically activating
protective mechanisms when anomalous observations are encountered.

\textbf{Limitations.}
First, while evaluation on two datasets (COHFACE and MAHNOB-HCI)
provides broader evidence than a single-dataset study, both involve
seated subjects in indoor settings.
Validation on datasets with more extreme non-stationary respiratory
dynamics (exercise recovery, sleep apnoea, neonatal breathing) is
needed to demonstrate the full potential of the adaptive mechanisms.
Second, while NAROSSM adapts $Q_t$ based on the estimated $R_t$,
the adaptation is a monotonic function of measurement noise;
more sophisticated joint $Q$--$R$ adaptation (e.g., dual estimation)
could improve tracking of rapidly changing respiratory dynamics.
Third, the ablation study reveals that robust z-score preprocessing
and adaptive filtering are complementary: removing clipping entirely
degrades NAROSSM's performance (OF-Farneb\"ack MAE increases from
0.280 to 0.355~BPM), while the Gaussian KF with fixed~$R$ is
unaffected.
This suggests that extreme outliers destabilise the adaptive
$R_t$ estimator through innovation variance inflation.
Future work should investigate more robust innovation statistics
(e.g., median-based rather than EMA-based) that are resilient to
extreme outliers without requiring preprocessing.

\textbf{Clinical significance.}
Beyond rate accuracy, NAROSSM's per-frame diagnostic outputs provide
information directly relevant to clinical deployment.
The $\lambdat$ trace identifies specific time intervals where the
observation signal is unreliable; the adaptive $\nuadapt$ provides
a continuous, per-family measure of observation quality.
The self-calibrating $\NIS \approx 1$ property ensures that
diagnostic thresholds have consistent statistical interpretation
across observation families and recording conditions, without
requiring family-specific tuning.


% ===================================================================
%  METHODS
% ===================================================================
\section*{Methods}

\subsection*{Problem formulation}

We model respiration as a latent quasi-periodic oscillatory process
observed through a noisy measurement channel.
The goal is to reconstruct the continuous respiratory trajectory
$x_1(t)$---representing chest-wall displacement---from a sequence
of scalar observations $\{y_1, y_2, \ldots, y_T\}$ obtained from
video-derived chest motion proxies.

\subsection*{Oscillatory state-space model}

The respiratory oscillation is represented by a two-dimensional state
vector $\state_t = [x_{1,t},\; x_{2,t}]^\top$, where $x_1$ and
$x_2$ are the in-phase and quadrature components of the oscillator.
Given an instantaneous frequency $f_t$ (adapted externally; see
below), the state transition is:
\begin{equation}
  \state_t = \mathbf{F}(f_t)\,\state_{t-1} + \mathbf{w}_t, \qquad
  \mathbf{F}(f_t) = \rho \begin{bmatrix}
    \cos\omega_t & \sin\omega_t \\
    -\sin\omega_t & \cos\omega_t
  \end{bmatrix},
  \label{eq:state_transition}
\end{equation}
where $\omega_t = 2\pi f_t \dt$ is the angular increment per sample,
$\rho = \exp(-1/f_s\tau_{\mathrm{env}})$ is an envelope damping factor
that prevents unbounded amplitude growth
($\tau_{\mathrm{env}}$ is the envelope time constant in seconds),
and $\mathbf{w}_t \sim \mathcal{N}(\mathbf{0}, \mathbf{Q})$ is
zero-mean Gaussian process noise with covariance
$\mathbf{Q} = \mathrm{diag}(q_x, q_x)$.

The observation equation is:
\begin{equation}
  y_t = \mathbf{H}\,\state_t + v_t, \qquad
  \mathbf{H} = [1,\; 0],
  \label{eq:observation}
\end{equation}
where $v_t$ is the measurement noise whose distribution is
\emph{not} assumed to be Gaussian but is instead modelled adaptively
as described below.

\textbf{Remark on linearity.}
For any given frequency $f_t$, the system
\eqref{eq:state_transition}--\eqref{eq:observation} is exactly
linear.
This is a deliberate design choice: by separating frequency
adaptation from state estimation, we avoid the nonlinear coupling
that would arise from including frequency as a third state variable,
and we can apply the standard Kalman filter without linearisation
approximations.

\subsection*{Harmonic-aware frequency initialisation}

The initial frequency $f_0$ is critical: the oscillatory state-space
model's narrow-band structure means that initialising at the wrong
frequency can cause persistent tracking error that the RTS smoother
cannot correct.
A common failure mode in respiratory signal processing is
\emph{harmonic confusion}: for slow breathing (below approximately
10~BPM), the chest-wall motion waveform is non-sinusoidal, causing
the second harmonic to dominate the Welch power spectral density.
A na\"ive spectral peak detector would lock onto the second harmonic
(e.g., 0.24~Hz instead of the true fundamental at 0.12~Hz).

We address this through a three-step harmonic refinement:
\begin{enumerate}[nosep]
\item \emph{Candidate detection.}
  A Welch periodogram with fine frequency resolution
  (60-second window at the signal's sampling rate) identifies the
  dominant spectral peak $f_{\mathrm{dom}}$ within the
  physiological respiratory band.
\item \emph{Subharmonic check.}
  The power at $f_{\mathrm{dom}}/2$ is compared to the dominant
  peak power.
  If the subharmonic has at least 15\% of the dominant peak's power,
  harmonic confusion is suspected.
\item \emph{Autocorrelation tiebreaker.}
  The autocorrelation function naturally identifies the fundamental
  period (not harmonics).
  If the autocorrelation peak frequency is closer to the subharmonic
  than to the dominant peak, the subharmonic is accepted as the
  true fundamental frequency.
\end{enumerate}

This conservative approach (requiring both spectral and
autocorrelation evidence) avoids false corrections while reliably
identifying harmonic confusion.
On COHFACE, harmonic refinement corrects the initial frequency on
approximately 8\% of trials, preventing systematic frequency
doubling errors.


\subsection*{Frequency adaptation module}

The instantaneous frequency $f_t$ is estimated by a hybrid module
that combines two complementary sources:

\textbf{Spectral source.}
Every $N_{\mathrm{spec}}$ frames, a Welch periodogram~\cite{Welch1967}
is computed over a sliding window of the most recent $W_{\mathrm{spec}}$
observations.
The spectral peak frequency $f_{\mathrm{spec}}$ within the
physiological respiratory band ($[f_{\min}, f_{\max}]$) provides a
stable but slowly adapting frequency estimate.

\textbf{Phase source.}
The instantaneous phase is computed from the state estimate:
$\phi_t = \mathrm{atan2}(x_{2,t}, x_{1,t})$.
The phase-based instantaneous frequency is:
\begin{equation}
  f_{\mathrm{phase},t} = \frac{f_s}{2\pi}
  \frac{d}{dt}\,\mathrm{unwrap}(\phi_t),
  \label{eq:phase_freq}
\end{equation}
which responds rapidly to frequency changes but is noise-sensitive.

\textbf{Self-gating combination.}
A critical design challenge is that frequency adaptation on
stable-breathing data can inject noise: the short-window spectral
estimate is inherently noisier than the initial full-signal estimate,
so adapting frequency when breathing rate is constant only degrades
the oscillator model.
We address this through a \emph{confirmation-gated} adaptation
mechanism: the spectral estimate must consistently disagree with
the current frequency for multiple consecutive update intervals
before adaptation occurs.

Formally, let $\delta_t = |f_{\mathrm{spec}} - f_{t-1}|$ be the
spectral deviation. A confirmation counter $c_t$ is maintained:
\begin{equation}
  c_t =
  \begin{cases}
    \max(c_{t-1} - 1,\; 0) & \text{if } \delta_t < \delta_{\min}, \\
    c_{t-1} + 1             & \text{otherwise},
  \end{cases}
  \label{eq:freq_confirm}
\end{equation}
where $\delta_{\min}$ is the minimum spectral deviation threshold.
Frequency adaptation occurs only when $c_t \geq c_{\mathrm{thresh}}$:
\begin{equation}
  f_t =
  \begin{cases}
    \mathrm{clip}\!\big(\alpha_f f_{\mathrm{spec}} + (1{-}\alpha_f) f_{t-1},\;
      f_{t-1} \pm \Delta f_{\max}\big) & \text{if } c_t \geq c_{\mathrm{thresh}}, \\
    f_{t-1} & \text{otherwise},
  \end{cases}
  \label{eq:freq_adapt}
\end{equation}
where $\alpha_f = \mathrm{clip}(\mathrm{SNR}/(\mathrm{SNR} + \beta), 0.05, 0.8)$
is the SNR-weighted blend coefficient and $\Delta f_{\max}$ is a
rate limiter.

This design is \emph{self-regulating}: on stable-breathing data
(e.g., COHFACE, where 64\% of trials exhibit zero frequency variation),
the spectral peak remains near $f_{t-1}$, the confirmation counter
stays low, and no adaptation occurs---the oscillator behaves
identically to a fixed-frequency baseline.
On data with genuine breathing rate changes, the persistent spectral
shift accumulates evidence in $c_t$ until the threshold is reached,
enabling smooth tracking.


\subsection*{Innovation-based adaptive measurement noise}

The measurement noise variance $R_t$ is estimated online from the
innovation sequence, following the principle of Mehra~\cite{Mehra1970}
and Mohamed and Schwarz~\cite{MohamedSchwarz1999}.

At each time step, the innovation is:
\begin{equation}
  e_t = y_t - \mathbf{H}\,\state_{t|t-1},
  \label{eq:innovation}
\end{equation}
where $\state_{t|t-1}$ is the predicted state.
The adaptive measurement noise estimate is updated via exponential
moving average:
\begin{equation}
  \hat{C}_{e,t} = \alpha_R\,\hat{C}_{e,t-1}
    + (1 - \alpha_R)\,e_t^2,
  \label{eq:ema_innovation}
\end{equation}
\begin{equation}
  R_t = \max\!\Big(\hat{C}_{e,t}
    - \mathbf{H}\,\mathbf{P}_{t|t-1}\,\mathbf{H}^\top,\;
    R_{\min}\Big),
  \label{eq:adaptive_R}
\end{equation}
where $\alpha_R \in (0, 1)$ is the forgetting factor and
$\mathbf{P}_{t|t-1}$ is the predicted state covariance.
The floor $R_{\min} = \alpha_{\mathrm{anchor}} \cdot R_0$ is anchored
to the initial measurement noise $R_0$, which prevents a positive
feedback instability identified during development: without anchoring,
overestimated $Q$ leads to $R_t \to 0$, which further suppresses
innovation variance, creating a self-reinforcing collapse.

\textbf{Empirical NIS.}
Rather than using the model-based NIS
($\NIS^{\mathrm{model}}_t = e_t^2 / S_t$, which is structurally
$\ll 1$ when $Q$ is overestimated), we compute the \emph{empirical}
NIS from the innovation autocovariance:
\begin{equation}
  \NIS^{\mathrm{emp}}_t = \frac{e_t^2}{\hat{C}_{e,t}},
  \label{eq:empirical_nis}
\end{equation}
which has expectation~1 by construction of the EMA estimator
$\hat{C}_{e,t}$.
This empirical NIS is used for all subsequent operations (Student-$t$
weight, kurtosis estimation, diagnostic output), ensuring that
$\NIS \approx 1$ regardless of the $Q$/$R$ balance.


\subsection*{Adaptive process noise (R-based bandwidth control)}

The process noise covariance $\mathbf{Q}_t$ controls the filter's
bandwidth: larger $Q_t$ allows the state to change more rapidly between
time steps (faster tracking but more noise), while smaller $Q_t$
constrains the state to follow the oscillatory model more closely
(smoother trajectories but slower adaptation).
Rather than fixing $Q_t$ as a constant, we use the adaptively
estimated $R_t$ as a per-frame signal quality indicator:
\begin{equation}
  q_{x,t} = q_x \cdot \mathrm{clip}\!\left(
    \left(\frac{R_t}{R_0}\right)^{\!\gamma},\;
    s_{\min},\; s_{\max}
  \right),
  \label{eq:adaptive_Q}
\end{equation}
where $q_x$ is the base process noise, $\gamma$ controls the
sensitivity (we use $\gamma = 0.5$, a square-root relationship),
and $[s_{\min}, s_{\max}]$ clips the scale factor to prevent
extreme values.

This creates a symmetric, signal-quality-dependent bandwidth control:
\begin{itemize}[nosep]
\item When $R_t < R_0$ (clean segment), the scale factor is $< 1$,
  reducing $Q_t$ and tightening the filter bandwidth.
  This \emph{anchors} the state estimate during clean windows,
  preventing the oscillator from drifting.
\item When $R_t > R_0$ (noisy segment), the scale factor is $> 1$,
  increasing $Q_t$ and widening the bandwidth to allow faster
  frequency tracking.
\end{itemize}

The clean-segment anchoring is critical for noisy families such as
P1D-DoF, where brief clean windows provide the most reliable
frequency information.
By reducing $Q_t$ during these windows, the filter ``locks on'' to
the correct frequency, improving robustness to subsequent noisy
segments.
On clean families (P1D-Quadratic), $R_t$ converges to a low value
throughout the recording, keeping $Q_t$ small and preserving the
spectral accuracy of the raw signal.


\subsection*{Kurtosis-adaptive Student-$t$ degrees of freedom}

The degrees-of-freedom parameter $\nu$ of the Student-$t$ observation
model controls the tail heaviness of the assumed noise distribution.
Rather than fixing $\nu$ as a hyperparameter, we adapt it from the
empirical kurtosis of the innovation sequence.

A running estimate of excess kurtosis is maintained:
\begin{equation}
  \hat{m}_{4,t} = \alpha_\kappa\,\hat{m}_{4,t-1}
    + (1-\alpha_\kappa)\,\tilde{e}_t^4,
  \qquad
  \hat{m}_{2,t} = \alpha_\kappa\,\hat{m}_{2,t-1}
    + (1-\alpha_\kappa)\,\tilde{e}_t^2,
  \label{eq:running_moments}
\end{equation}
\begin{equation}
  \kurtosis = \frac{\hat{m}_{4,t}}{\hat{m}_{2,t}^2} - 3,
  \label{eq:running_kurtosis}
\end{equation}
where $\tilde{e}_t = e_t / \sqrt{S_t}$ is the standardised
innovation.
The adaptive degrees of freedom are obtained by inverting the
Student-$t$ excess kurtosis relation:
\begin{equation}
  \nuadapt = \mathrm{clip}\!\left(
    \frac{6}{\max(\kurtosis,\, \epsilon)} + 4,
    \;\nu_{\min},\; \nu_{\max}
  \right),
  \label{eq:adaptive_nu}
\end{equation}
where $\epsilon > 0$ prevents division by zero, and
$[\nu_{\min}, \nu_{\max}]$ bounds the parameter to a numerically
stable range (e.g., $[3, 200]$).

When innovations are heavy-tailed ($\kurtosis \gg 0$), $\nuadapt$
decreases, producing aggressive outlier suppression.
When innovations are near-Gaussian ($\kurtosis \approx 0$),
$\nuadapt$ increases toward $\nu_{\max}$, and the filter converges
to standard Gaussian Kalman filtering.
This constitutes a direct, data-driven implementation of
inference-to-noise alignment.


\subsection*{Student-$t$ robust Kalman update}

Given the adaptively estimated $R_t$ and $\nuadapt$, the measurement
update follows the variational Bayes
approach~\cite{Huang2017,Agamennoni2012}.
The key quantity is the scale weight:
\begin{equation}
  \lambdat = \frac{\nuadapt + 1}{\nuadapt + \NIS^{\mathrm{emp}}_t},
  \label{eq:lambda}
\end{equation}
which satisfies $0 < \lambdat \le 1$ for finite $\nuadapt$ when
$\NIS^{\mathrm{emp}}_t \geq 1$.
When the empirical NIS is near unity (normal observation),
$\lambdat \approx 1$ and the update is nearly Gaussian.
When the innovation is anomalously large,
$\lambdat \ll 1$ and the effective measurement noise
$\Reff = R_t / \lambdat$ increases, reducing the Kalman gain and
suppressing the influence of the outlier observation.

The Kalman gain is then:
\begin{equation}
  \mathbf{K}_t = \mathbf{P}_{t|t-1}\,\mathbf{H}^\top \,
  \big(\mathbf{H}\,\mathbf{P}_{t|t-1}\,\mathbf{H}^\top
    + \Reff\big)^{-1},
  \label{eq:kalman_gain}
\end{equation}
and the state and covariance updates follow the Joseph
stabilised form:
\begin{align}
  \state_{t|t} &= \state_{t|t-1}
    + g_t\,\mathbf{K}_t\, e_t,
    \label{eq:state_update} \\
  \mathbf{P}_{t|t} &= (\mathbf{I} - g_t\,\mathbf{K}_t\,\mathbf{H})
    \,\mathbf{P}_{t|t-1}\,
    (\mathbf{I} - g_t\,\mathbf{K}_t\,\mathbf{H})^\top
    + g_t^2\,\mathbf{K}_t\,\Reff\,\mathbf{K}_t^\top,
    \label{eq:cov_update}
\end{align}
where $g_t \in [0, 1]$ is the \emph{oscillator-native trust gate},
a key design element motivated by a challenge that deep learning
approaches handle implicitly through learned attention mechanisms
but that classical estimators must address through explicit design.

The question of \emph{which observations to trust} is fundamental:
a standard Kalman filter treats all observations equally, while
learned models can automatically identify unreliable inputs through
data-driven attention weights.
We address this challenge by deriving $g_t$ from signals that are
\emph{intrinsic to the oscillatory state-space model itself}---not
from the innovation magnitude (NIS), which creates a bistable
feedback loop (high NIS $\to$ gate closes $\to$ filter diverges
$\to$ NIS remains high).

Specifically, $g_t$ combines two oscillator-native quality measures
with a threshold mechanism that ensures clean observations receive
full trust:
\begin{equation}
  g_t =
  \begin{cases}
    1 & \text{if } q_{\phi,t} \cdot q_{A,t} \geq \tau_g, \\
    \max\!\big(q_{\phi,t} \cdot q_{A,t},\; g_{\min}\big) & \text{otherwise},
  \end{cases}
  \label{eq:gate}
\end{equation}
where $\tau_g$ is a trust threshold and $g_{\min} > 0$ prevents
complete observation rejection.
The threshold ensures that normal physiological variations in
amplitude and phase do not trigger unnecessary gain reduction,
so that NAROSSM preserves the full Kalman gain on clean frames
(behaving identically to a standard filter) while activating
proportional gating only on genuine outlier frames.

\textbf{Phase coherence} ($q_{\phi,t}$) measures whether the
observation is consistent with the oscillator's predicted phase
progression:
\begin{equation}
  q_{\phi,t} = \exp\!\left(-\frac{(\Delta\phi_t -
    \omega_t \dt)^2}{2\sigma_\phi^2}\right),
  \label{eq:phase_coherence}
\end{equation}
where $\Delta\phi_t = \phi(\hat{\state}_{t|t}) - \phi(\state_{t-1|t-1})$
is the actual phase change and $\omega_t \dt = 2\pi f_t \dt$ is the
expected phase increment.
A normal breathing observation produces smooth phase progression
($q_{\phi,t} \approx 1$); a motion artefact causes an abrupt phase
jump ($q_{\phi,t} \ll 1$).

\textbf{Amplitude stability} ($q_{A,t}$) detects sudden changes in
the oscillator envelope:
\begin{equation}
  q_{A,t} = \exp\!\left(-\frac{1}{2}\left(
    \frac{|A_t - \bar{A}_t|}{\bar{A}_t \cdot \sigma_A}
    \right)^2\right),
  \label{eq:amp_stability}
\end{equation}
where $A_t = \|\state_t\|$ is the instantaneous amplitude and
$\bar{A}_t$ is an exponentially weighted moving average.
Respiratory amplitude varies slowly; abrupt jumps indicate artefacts.

This gate design is \emph{self-stabilising}: when $g_t$ reduces the
observation influence, the filter relies on its internal oscillatory
model, which naturally produces smooth phase and stable amplitude.
This restores $q_{\phi,t}$ and $q_{A,t}$ toward unity, reopening the
gate---in contrast to NIS-based gates, which lock into a closed state.
The trust gate provides a \emph{prior} assessment of observation
reliability (model-prediction consistency), complementing the
Student-$t$ weight $\lambdat$ which provides a \emph{posterior}
assessment (innovation statistics).


\subsection*{Rauch--Tung--Striebel backward smoother}

After the forward filtering pass, a backward
Rauch--Tung--Striebel (RTS) smoothing
pass~\cite{RauchTungStriebel1965,Sarkka2013} refines the state
estimates using future observations:
\begin{align}
  \mathbf{G}_t &= \mathbf{P}_{t|t}\,\mathbf{F}_{t+1}^\top\,
    \mathbf{P}_{t+1|t}^{-1},
    \label{eq:rts_gain} \\
  \state_{t|T} &= \state_{t|t} + \mathbf{G}_t\,
    (\state_{t+1|T} - \state_{t+1|t}),
    \label{eq:rts_state} \\
  \mathbf{P}_{t|T} &= \mathbf{P}_{t|t} + \mathbf{G}_t\,
    (\mathbf{P}_{t+1|T} - \mathbf{P}_{t+1|t})\,
    \mathbf{G}_t^\top.
    \label{eq:rts_cov}
\end{align}

Because the state-space model is exactly linear for any given
frequency sequence, the RTS smoother is the optimal (minimum
variance) fixed-interval smoother for this system.
This contrasts with extended Kalman smoother approaches applied to
nonlinear state-space models, which inherit the linearisation
errors of the forward EKF pass.

The smoothed state $\state_{t|T}$ provides the final respiratory
trajectory estimate, with $x_{1,t|T}$ representing the reconstructed
respiratory waveform.


\subsection*{Algorithm summary}

Algorithm~\ref{alg:narossm} summarises the complete NAROSSM
forward pass.

\begin{algorithm}[!ht]
\caption{NAROSSM Forward Pass}
\label{alg:narossm}
\begin{algorithmic}[1]
\Require Observations $\{y_1, \ldots, y_T\}$, initial state
  $\state_0$, $\mathbf{P}_0$, parameters
  $\alpha_R$, $\alpha_\kappa$, $\mathbf{Q}$, $R_{\min}$,
  $[\nu_{\min}, \nu_{\max}]$
\For{$t = 1, \ldots, T$}
  \State \textbf{Frequency adaptation:}
    Update $f_t$ via self-gating spectral module
    \eqref{eq:freq_confirm}--\eqref{eq:freq_adapt}
  \State \textbf{Predict:}
    $\state_{t|t-1} = \mathbf{F}(f_t)\,\state_{t-1|t-1}$;\quad
    $\mathbf{P}_{t|t-1} = \rho^2\,\mathbf{R}(\omega_t)\,
      \mathbf{P}_{t-1|t-1}\,\mathbf{R}(\omega_t)^\top + \mathbf{Q}$
  \State \textbf{Innovation:}
    $e_t = y_t - \mathbf{H}\,\state_{t|t-1}$
  \State \textbf{Adaptive $R_t$:}
    Update via \eqref{eq:ema_innovation}--\eqref{eq:adaptive_R}
  \State \textbf{Adaptive $Q_t$:}
    $q_{x,t} = q_x \cdot \mathrm{clip}((R_t/R_0)^\gamma,\, s_{\min},\, s_{\max})$
    via \eqref{eq:adaptive_Q}
  \State \textbf{Innovation variance:}
    $S_t = \mathbf{H}\,\mathbf{P}_{t|t-1}\,\mathbf{H}^\top + R_t$
  \State \textbf{Adaptive $\nu_t$:}
    Update running kurtosis and $\nuadapt$ via
    \eqref{eq:running_moments}--\eqref{eq:adaptive_nu}
  \State \textbf{Student-$t$ weight:}
    $\lambdat = (\nuadapt + 1) / (\nuadapt + e_t^2/S_t)$
  \State \textbf{Effective noise:}
    $\Reff = R_t / \lambdat$
  \State \textbf{Kalman gain:}
    $\mathbf{K}_t = \mathbf{P}_{t|t-1}\,\mathbf{H}^\top\,
      (S_t')^{-1}$
    where $S_t' = \mathbf{H}\,\mathbf{P}_{t|t-1}\,\mathbf{H}^\top
      + \Reff$
  \State \textbf{Trust gate:}
    Compute $q_{\mathrm{raw}} = q_{\phi,t} \cdot q_{A,t}$;
    \textbf{if} $q_{\mathrm{raw}} \geq \tau_g$ \textbf{then} $g_t = 1$
    \textbf{else} $g_t = \max(q_{\mathrm{raw}}, g_{\min})$
    via \eqref{eq:gate}
  \State \textbf{Gated update:}
    $\state_{t|t}$, $\mathbf{P}_{t|t}$ via
    \eqref{eq:state_update}--\eqref{eq:cov_update} with $g_t$
  \State \textbf{Diagnostics:}
    Record $\NIS_t$, $\lambdat$, $R_t$, $\nuadapt$, $g_t$
\EndFor
\State \textbf{RTS smoother:} Backward pass
  \eqref{eq:rts_gain}--\eqref{eq:rts_cov}
\State \Return $\{\state_{t|T}\}_{t=1}^T$ and diagnostic traces
\end{algorithmic}
\end{algorithm}

% --- Figure 6: Algorithm block diagram ---
\begin{figure}[!ht]
\centering
% \includegraphics[width=\linewidth]{figures/fig6_algorithm_block}
\fbox{\parbox{0.95\linewidth}{\centering\vspace{4cm}
\textbf{Figure 6: NAROSSM Algorithm Block Diagram}\\[6pt]
Signal flow: $y_t \to$ Innovation $e_t \to$ Adaptive $R_t$ $\to$\\
Running kurtosis $\to$ Adaptive $\nu_t$ $\to$ Student-$t$ $\lambda_t$ $\to$\\
$R_{\mathrm{eff}}$ $\to$ Kalman Gain $\to$ State Update $\to$ RTS\\[6pt]
\emph{(Detailed block diagram to be created in PaperBanana)}
\vspace{4cm}}}
\caption{%
  Block diagram of the NAROSSM algorithm showing the information
  flow between the adaptive noise estimation, Student-$t$ robust
  update, and state estimation components.
  Solid arrows indicate the main signal path; dashed arrows indicate
  diagnostic outputs.
}
\label{fig:algorithm_block}
\end{figure}


\subsection*{Observation extraction}

Five observation families are employed, each representing a different
approach to extracting chest-wall motion from video:

\begin{itemize}[nosep]
\item \textbf{OF-Farneb\"ack}:
  Dense optical flow~\cite{Farneback2003} computed over the
  upper-body region of interest (ROI), with the vertical component
  averaged to produce a scalar respiratory proxy.
\item \textbf{Profile1D-Linear, -Quadratic, -Cubic}:
  One-dimensional spatial intensity profiles extracted along vertical
  lines within the chest ROI, fitted with polynomials of degree~1,
  2, and~3 respectively.
  The temporal variation of profile coefficients captures
  chest-wall displacement.
\item \textbf{Profile1D-DoF}:
  Degrees-of-freedom-based profile extraction that adapts the
  polynomial order to the local signal structure.
\end{itemize}


\subsection*{Datasets}

\textbf{COHFACE}~\cite{COHFACE2017} contains 160~video recordings from
40~subjects (4~sessions each) in a controlled laboratory environment.
Videos are recorded at 20~fps ($640 \times 480$) with simultaneous
ground-truth respiratory signals from a chest-mounted BVP sensor.
Each recording is approximately 60~seconds.
The controlled protocol (seated, still, neutral task) yields stable
breathing patterns.

\textbf{MAHNOB-HCI}~\cite{MAHNOB2012} is an emotion-elicitation
dataset containing 525~recordings from 27~subjects.
Videos are recorded at 61~fps ($780 \times 580$) with ground-truth
respiration from a chest belt (BDF channel~44, 256~Hz).
Each recording is approximately 120~seconds.
The emotion-elicitation protocol (subjects watching affective video
clips) induces more naturalistic and variable breathing patterns,
including sighs, irregular rhythms, and body movements.

Both datasets use upper-body (chest) ROI extraction for all five
observation families.
Preprocessing, state-space model configuration, and evaluation
metrics are identical across datasets to enable direct comparison.


\subsection*{Evaluation protocol}

\textbf{Frequency-domain metrics.}
Respiratory rate is estimated from the reconstructed signal
$\hat{x}_{1,t|T}$ using Welch PSD~\cite{Welch1967} peak detection
over 30-second sliding windows with 1-second stride.
Five complementary metrics are reported:
(1)~windowed mean absolute error (MAE, BPM);
(2)~root mean squared error (RMSE, BPM);
(3)~Pearson correlation coefficient ($r$) between estimated and
ground-truth rate trajectories;
(4)~spectral signal-to-noise ratio (SNR\textsubscript{spec}, dB),
computed as the ratio of the respiratory peak power to the
out-of-band spectral floor; and
(5)~Kullback--Leibler divergence between estimated and ground-truth
normalised power spectral densities, measuring overall spectral
shape agreement.

\textbf{Time-domain metrics.}
Waveform reconstruction fidelity is assessed on bandpass-filtered,
cross-correlation-aligned signals using six metrics:
(1)~Lin's concordance correlation coefficient
(CCC)~\cite{LinCCC1989}, which jointly measures precision and
accuracy of waveform agreement;
(2)~mean absolute error (MAE) and (3)~root mean squared error (RMSE)
on aligned waveforms;
(4)~time-domain signal-to-noise ratio (SNR\textsubscript{time}, dB);
(5)~dynamic time warping (DTW) distance, which captures waveform
shape similarity robust to temporal misalignment; and
(6)~peak-to-peak interval (PPI) MAE, measuring the accuracy of
individual breath cycle duration estimation.

\textbf{Statistical testing.}
All comparisons use paired Wilcoxon signed-rank tests with Bonferroni
correction for multiple comparisons.
A result is considered statistically significant at $p < 0.05$
(corrected).


\subsection*{Baselines}

Two baselines are compared, chosen to isolate the contribution of
NAROSSM's adaptive mechanisms from the underlying state-space
structure:
\begin{itemize}[nosep]
\item \textbf{Base (raw spectral):}
  Direct Welch PSD peak detection on the preprocessed observation
  signal without any temporal filtering.
  This establishes the performance achievable from spectral analysis
  alone.
\item \textbf{Gaussian KF (kfstd):}
  The \emph{same} two-dimensional oscillatory state-space model
  (Equation~\ref{eq:state_transition}--\ref{eq:observation}) with a
  standard Kalman filter (fixed $R$, Gaussian $v_t$), fixed
  transition frequency $f$ estimated from the initial spectral peak,
  and RTS backward smoother.
  This baseline shares identical preprocessing, state-space
  dimensionality, base process noise $Q$, and smoothing algorithm with
  NAROSSM; the \emph{only} differences are NAROSSM's adaptive $R_t$,
  R-based adaptive $Q_t$, kurtosis-adaptive $\nuadapt$, Student-$t$
  update, oscillator-native trust gate, and harmonic-aware frequency
  initialisation.
  This controlled comparison ensures that any performance difference
  is attributable solely to the noise-adaptive mechanisms.
\end{itemize}


% ===================================================================
%  ACKNOWLEDGEMENTS, CONTRIBUTIONS, ETC.
% ===================================================================

\bibliography{refs}

\section*{Acknowledgements}

This work was supported by [funding information].

\section*{Author contributions statement}

Y.-J.C.\ conceived the research framework, designed the NAROSSM
algorithm, implemented the software, conducted all experiments, and
wrote the manuscript.
D.-Y.K.\ supervised the research, provided critical feedback on the
methodology, and reviewed the manuscript.

\section*{Additional information}

\textbf{Data availability.}
The COHFACE dataset is publicly available at
\url{https://www.idiap.ch/dataset/cohface}.
The MAHNOB-HCI dataset is available at
\url{https://mahnob-db.eu/hci-tagging/}.
Code implementing NAROSSM and reproducing all experiments will be
made available upon publication.

\textbf{Competing interests.}
The authors declare no competing interests.


\end{document}
